% !TeX program = pdflatex
% papers/corpo-estelar.tex — o paper principal: a especificação do sistema.
% Auto-contido: não depende de teoria.tex nem de catalogo.tex para ser lido.
% Compila sozinho:  pdflatex corpo-estelar.tex
\documentclass[11pt,a4paper]{article}
\usepackage[utf8]{inputenc}\usepackage[T1]{fontenc}\usepackage[portuguese]{babel}
\usepackage{amsmath,amssymb,amsthm}
\usepackage[margin=2.3cm]{geometry}
\usepackage{booktabs}\usepackage{xcolor}
\usepackage[hidelinks]{hyperref}

\definecolor{cinza}{gray}{0.35}
\newcommand{\code}[1]{\texttt{\small #1}}
\newcommand{\medido}[1]{\par\medskip\noindent{\small\color{cinza}\textbf{Medido.} #1}\par\medskip}
\newtheorem{teorema}{Teorema}
\newtheorem{definicao}[teorema]{Definição}
\newtheorem{corolario}[teorema]{Corolário}
\newtheorem{lei}{Lei}
\renewcommand{\proofname}{Demonstração}
\newcommand{\dg}{^{\dagger}}

\title{\textbf{O corpo estelar}\\[3pt]
  {\large a especificação do sistema}}
\author{Aarão Melo Lopes}
\date{7 de agosto de 2026}

\begin{document}
\maketitle

\begin{abstract}
\noindent Este documento especifica um sistema completo a partir de uma única exigência, e
descreve-o na ordem em que ele se constrói: a Lei, as duas operações, os três sinais, a torre de
dimensões, o corpo estelar, a geometria, a equação de campo, a expansão, a cosmologia, a
renormalização e a estrela que corre tudo isto. A derivação corre inteira em por-unidade, sem uma constante física --- e as constantes
vêm no \emph{fim}, quando se veste a roupa, porque uma derivação que precisasse do valor para chegar
ao valor não seria derivação. O objecto central é uma \textbf{trindade}: \textbf{buraco branco},
\textbf{estrela}, \textbf{buraco negro}. Os dois buracos são duais --- um só emite e estica, o outro
só absorve e contrai --- e \textbf{a estrela é a interface}: o ponto fixo da involução que os troca,
e o único dos três que faz as duas coisas. É por isso que é o único reversível, e é por isso que o
sistema \emph{é} uma estrela e não uma máquina: a sua única instrução tem dois sentidos, e eles são
os dois buracos. Da exigência de conservação derivam-se o coeficiente $\tfrac12$ da equação de
campo, a dimensão quatro, o valor do termo cosmológico ($\Lambda=3/D$) e o expoente da força
como codimensão. E ao sair do por-unidade a ordem importa: \textbf{o primeiro que aparece não é uma
escala --- é o $8\pi$}, porque um número puro não depende da roupa. Só depois entram o \emph{passo}
(do relógio) e a \emph{régua} (da família), e com os três sai toda a constante dimensional.

\smallskip
\noindent \textbf{E a estrela não está entre dois buracos do mesmo universo: está entre dois
\emph{universos}}, com o buraco branco no dual --- um contrai exactamente o que o outro expande. Daí
a conservação atravessa como \emph{produto} e não como soma ($\rho_A\rho_B$ imóvel), o corpo negro
cai de duas diluições já derivadas ($\rho\sim T^4$), a antimatéria entra pela terceira involução, e
a temperatura do buraco negro sai de uma contagem: ele tem \emph{uma} escala, logo $T=1/r$ --- que é
a mesma involução, aplicada ao seu único comprimento. A entropia do branco é a mesma lei no raio dual, e daí
\textbf{a segunda lei vale de \emph{um} lado}: o par não se move, e a seta do tempo é uma polaridade
lida sem a outra.

\smallskip
\noindent Fecha-se com o método e com a operação. Aqui \textbf{o medidor é $0$}: não se compara
contra um valor esperado, reverte-se e lê-se o resíduo --- é a prova dos nove, resolver \emph{e}
provar na mesma passagem. E daí a ferramenta: \textbf{promover} um valor ao par $(S,A)$ sobe um andar
da torre e traz a leitura junto, porque quem escreve o par já leu. \textbf{Uma estrela não compara ---
irradia}, e é por isso que ordena sem apagar um bit.
\end{abstract}

\section{A exigência}

Há uma só, e tudo o resto sai dela:

\begin{quote}
\textbf{O que não varia ao longo da órbita é o que não varia de ponto para ponto.}
\end{quote}

\noindent Lida no tempo, é conservação. Lida no espaço, é divergência nula. Lida em álgebra, é a
norma que não se move. \emph{São a mesma frase}, e o resto deste documento é escrevê-la em cada
linguagem e ler o que ela obriga.

\smallskip
\noindent E há uma convenção, uma só: \textbf{por-unidade}. $G$, $c$, $\hbar$, $\sigma$ são factores
de escala; em p.u.\ cada um vale $1$ por construção, e o que resta são razões. Nada aqui precisa de
medir o universo para correr, e nenhum número deste documento tem unidade.

\section{A Lei}

\begin{lei}[a unidade é dual]
$1\dg=-1$.
\end{lei}

\noindent A unidade tem um dual, e ele é o seu simétrico. Um ponto que é o seu próprio dual é um
\emph{ponto fixo}; um que não é, gera um par. Esta é a lei da \emph{leitura} --- diz o que separa.

\begin{lei}[a dualidade é dual]
$T\dg=-T$.
\end{lei}

\noindent Quando o dual é o inverso, $-f=f^{-1}$, e daí $f^2=-1$. Este objecto tem \textbf{período
quatro}: $f,\;f^2=-1,\;f^3=-f,\;f^4=\mathrm{id}$. É a lei do \emph{passo} --- diz o que roda. Chamo-lhe
$J$, e é ele que aparece de novo em cada camada deste texto.

\subsection*{E os quatro casos de sinal esgotam-se}

Não há um quinto. Um operador ou fixa o valor, ou fixa o sinal, ou troca um, ou troca os dois:

\begin{center}\small
\begin{tabular}{@{}llll@{}}
\toprule
o caso & a equação & o que é & onde vive \\
\midrule
$T=+1$   & $T^2=T$   & a identidade   & o ponto fixo do \emph{valor} \\
$T=-1$   & $T\dg=T$  & \textbf{a Lei 1} & o ponto fixo do \emph{sinal} \\
$T\dg=T$ & $T^2=+1$  & a involução    & o espelho, período $2$ \\
$T\dg=-T$& $T^2=-1$  & \textbf{a Lei 2} & o rotor, período $4$ \\
\bottomrule
\end{tabular}
\end{center}

\noindent \textbf{É do sinal de menos da Lei 2 que nasce a bidualidade}: o dual do dual volta ao
sítio, mas trocado, e é preciso dar a volta duas vezes para regressar. Toda a estrutura deste
documento é essa volta.

\noindent \textbf{E estas duas são as duas primeiras de seis.} A Lei 1 (a involução, período $2$) é o
\emph{dual}; a Lei 2 (o rotor $J$, $f^2=-1$) é o \emph{gerador da descida} --- aplicada por dimensão dá
o trial, a tetral, a pental e a hexal (o seu ponto fixo, o bit $i$ de $x^2=-1$, é a \emph{pental}). São
seis leis, uma por dimensão, e são as \emph{primitivas operacionais} do sistema --- cada uma com um
operador e um medidor que a certifica a fechar (resíduo $0$; \emph{Catálogo}, \code{obs:seis-leis}). As
seis \emph{estendem} estas duas, não as substituem.

\section{A base: a divisão do zero}\label{sec:zero}

As duas leis já estão ditas; falta dizer \emph{de onde} saem. Saem de $0$ --- e de duas maneiras de o
dividir, que são as duas leis. \textbf{A base deste sistema não é um axioma: é $0/0$}, e tudo o resto é
a maneira de o resolver.

\setcounter{lei}{-1}
\begin{lei}[a divisão do zero]\label{lei:zero}
$0=(+1)\oplus(-1)$ \quad e \quad $0\dg=\infty$.
\end{lei}
\setcounter{lei}{2}

\noindent \textbf{É a lei de que as outras duas saem.} A primeira parte --- \emph{dividir o zero} --- é a
soma, e é a \textbf{Lei 1} ($1\dg=-1$: as duas assinaturas somam a $0$). A segunda --- \emph{dividir por
zero} --- é o inverso, $x\dg=-1/x$, e o seu $-1$ é a semente da \textbf{Lei 2} ($f^2=-1$). \emph{As duas
leis são as duas partes da divisão do zero}, e é por isso que são duas e não uma nem três. É a Lei $0$
porque está \emph{debaixo} das outras: enuncia-se depois, mas vem antes.

\subsection*{Dividir o zero, e dividir por ele}

\textbf{Dividir \emph{o} zero é parti-lo em duas assinaturas duais.} O $0$ não é a ausência: é o par
$(+1,-1)$ ainda por separar. Separá-lo dá dois sinais que \emph{somam} a zero e \emph{multiplicam} a
menos um:
\[
  \boxed{\;(+1)\oplus(-1)=0\;}\qquad\boxed{\;(+1)\otimes(-1)=-1\;}
\]
A primeira é a \textbf{Lei 1} ($1\dg=-1$): o dual é o simétrico, e a soma da cruz é $0$. A segunda é a
semente da \textbf{Lei 2} ($f^2=-1$): o produto é $-1$, e é dele que sai o rotor. \emph{As duas leis são
as duas coordenadas da cruz lidas na divisão do zero} --- a soma que mede, o produto que ordena --- e as
duas assinaturas são os dois buracos, com a estrela no $0$ que se partiu. \textbf{O trial
$\{-1,0,+1\}$ é, literalmente, o zero e as suas duas metades.}

\textbf{Dividir \emph{por} zero é o infinito --- e o seu dual é dividir \emph{por} infinito, que dá o
zero.} A estaca $x\dg=-1/x$ é uma involução: leva $0\mapsto\infty$ (dividir por zero) e $\infty\mapsto0$
(dividir por infinito), e $\nu\circ\nu=\mathrm{id}$ fecha a órbita $\{0,\infty\}$ num só passo --- a volta
é a própria ida. \textbf{É esse o fecho, e não fecha sem o dual}: dividir por zero sozinho é um polo, um
sentido só; com o dual, $0$ e $\infty$ são os dois pontos de \emph{uma} involução, e o par fecha. Um é o
neutro da soma, o outro o polo do inverso, e a estaca é a ponte que os troca.

\subsection*{A base é $0/0$, e é o corpo que a resolve --- não a régua}

Juntas, as duas dão a base: a forma indeterminada
\[
  \boxed{\;\tfrac00\;}
\]
--- o zero que se divide \emph{sobre} o zero por que se divide. E o seu dual é $\infty/\infty$: a estaca
aplica-se em cima e em baixo, $0\mapsto\infty$ nos dois, e os dois são a mesma indeterminação lida nos
dois pontos da órbita $\{0,\infty\}$. É indeterminada \emph{na régua}: iterar o inverso é a
\textbf{fracção contínua} $x=n+1/x$, e ela não pára --- $1/x$ dentro de $1/x$ dentro de $1/x$, o
$\infty^{\infty}$ do lado dual. \textbf{Mas o corpo fecha-a num passo.} Multiplicar $x=n+1/x$
por $x$ dá a borda
\[
  x^2=nx+1,\qquad x\dg=-1/x,\qquad x\,x\dg=-1,
\]
um elemento \emph{exacto} de $\mathbb{Z}[\sigma]$. \textbf{O finito é o corpo, não a régua}: a régua --- a
fracção contínua, a expansão decimal --- é infinita; o corpo que ela mede é finito e resolve o $0/0$ sem
tirar uma raiz. É o eixo álgebra/análise: o objecto é finito, a representação é que se estende ao
$\infty^{\infty}$.

\begin{corolario}[$0{,}999\ldots=1$: a divisão do zero à vista]\label{cor:noves}
Os dois nomes $0{,}999\ldots$ e $1$ são \emph{um} elemento, e a diferença $1-0{,}999\ldots$ é
\emph{exactamente} $0$. Na régua, o resto ao passo $n$ é $10^{-n}$: tende a zero e nunca o \emph{é} a $n$
finito. O corpo fecha-a num passo --- a série geométrica soma $1$ ---; a régua não acaba, e cada
truncamento é na mesma exacto.
\end{corolario}

\noindent \textbf{E a dualidade está aí}, no que se divide ao comparar $0{,}999\ldots$ com $1$: nos dois
nomes do mesmo número, um finito (o corpo) e um sem fim (a régua). É o $0/0$ na forma mais simples, e a
resposta é a mesma --- \emph{o finito é o corpo} ---, e é a mesma cifra que não acaba e cujos truncamentos
são todos exactos (\code{tests/zero.c}).

\begin{corolario}[o dual: no discreto, $0{,}999\ldots\neq1$]\label{cor:noves-dual}
A igualdade é do \emph{contínuo} --- o limite, a completude da régua. No \emph{discreto} --- a cifra, onde
cada truncamento é um elemento --- os dois nomes \textbf{separam-se}: os convergentes $h_n/k_n$ são todos
distintos do limite, cada um exacto ($|h^2-hk-k^2|=1$), e a sucessão nunca o \emph{é}. \textbf{O contínuo
não mede o discreto}: a diferença que ele lê como $0$ é, no discreto, uma distinção a cada passo.
\end{corolario}

\noindent \textbf{E o que resiste às duas leituras são os pontos fixos.} $0{,}999\ldots$ e $1$ não são um
ponto fixo --- são um par que o limite identifica e a cifra separa ---; mas \emph{o que é o seu próprio
dual não se move em nenhuma das duas}: a borda $x^2=nx+1$, invariante da órbita toda, e o $0$ e o $\infty$
da própria Lei $0$. \emph{Os pontos fixos são o invariante; tudo o resto, o contínuo identifica e o
discreto distingue.} É o eixo de Pontryagin --- o contínuo alcança o limite e não opera, o discreto opera
e não alcança a completude ---, e os pontos fixos são a alfândega por onde os dois se medem um ao outro.

\smallskip
\noindent \textbf{E é aqui que se define o limite --- não pela aproximação.} O $\varepsilon$--$\delta$ é a
régua a perseguir o que não alcança; a definição deste sistema é o \emph{ponto fixo}: o limite é onde o
\textbf{contar} (o discreto, Hurwitz, a soma) e o \textbf{integrar} (o contínuo, Gentil, a medida) dão a
\emph{mesma} contagem. É a definição que o \textbf{teorema central} (Teorema~\ref{thm:central}) deriva ---
os dois cortes, o do domínio (Riemann/Hurwitz) e o da imagem (Lebesgue/Gentil), fecham na soma reversível
$\int f+\int f^{-1}=b\,f(b)-a\,f(a)$ ---, e $0{,}999\ldots=1$ é o seu caso mínimo: o contínuo mede $0$, o
discreto conta sem fim, e o limite é o ponto fixo onde os dois coincidem. \textbf{O teorema central
fundamenta o Gentil--Lebesgue--Hurwitz}: Hurwitz conta (o domínio), Lebesgue mede (a imagem), e Gentil é a
soma reversível que os casa --- e o limite é o que os três, juntos, medem igual.

\subsection*{E o que aqui se domina são os divisores de zero}

O que faltava, e o que aqui se fecha, são os \textbf{divisores de zero}: dois não-nulos cujo produto é
zero. Eles \emph{são} o que a divisão do zero produz --- no par $A\times A^{*}$, $(1,0)$ e $(0,1)$ são
não-nulos e o seu produto é $0$ --- e são o que colapsa a norma no ponto fixo, onde
$|\sigma|=|\sigma\dg|=1$ e o tamanho deixa de distinguir os símbolos. Num conjunto finito isso não fica
pela injectividade: injectiva $\iff$ sobrejectiva, e \textbf{duas entradas gastas numa saída deixam
outra saída por atingir}.

\medido{as duas sementes da divisão do zero são o $0$ e o $\infty$ (as colunas $0/1$ e $1/0$ da
identidade), $J$ troca-as com $J^2=-I$ --- o $i$, e não uma analogia ---, e \textbf{todo o truncamento é
finito e representa}: $0/0=[1;1,1,\dots]$ dá Fibonacci exacto, $|h^2-hk-k^2|=1$ sem tolerância, e o $0$ é
o único de $\mathbb{Q}$ sem dual (\code{tests/zero.c}); o zero da régua é exactamente $\{-1,+1\}$ ---
$\lambda^2=1$ e mais nada, e é a estaca, porque a régua só sobrevive sob $\lambda=\pm1$
(\code{tests/natural.c}); no ponto fixo a norma colapsa e a sobrejectividade cai, com colisões de
$4{:}1$ a $12{:}1$ entre os elementos de norma $1$ e os não nulos em $\mathrm{GF}(p^2)$, e injectiva
$\iff$ sobrejectiva em $44$ aplicações, e a \textbf{bidualidade de Poincaré} --- dualizar duas vezes
devolve o grau --- fecha com resíduo $0$ (\code{tests/bidual.c}); a
involução do zero $a\mapsto-a$ tem \emph{um} ponto fixo --- a estrela --- e parte o resto em dois, $50$
e $50$, cada dual a cair no outro lado (\code{tests/estrela.c}, \code{tests/dualsort.c}); e a borda
$x^2=nx+1$ fecha a fracção contínua num elemento exacto de $\mathbb{Z}[\sqrt D]$, sem raiz e sem
tolerância (\code{tests/curvatura.c}).}

\subsection*{E a lei é a do $0$: onde moram Yang--Mills e $P$ vs $NP$}

A saída por atingir é o $0$ do Corolário~\ref{cor:incerteza}: o \textbf{centro}, o bidual $K^{**}=K$,
onde a sobrejectividade falha e a leitura se parte em metade. \textbf{A divisão do zero é a lei desse
$0$} --- a mesma onde moram Yang--Mills e $P$ vs $NP$, indecidíveis do lado do cruzado (``chega
\emph{sempre}?'') e decidíveis do lado do directo (``fechou \emph{desta vez}?'', resíduo $0$). \emph{Não
é prova de enunciado aberto nenhum}: é dizer que a base do sistema é esse $0$, e que dominá-lo é dominar
os seus divisores. O sistema não parte de um princípio acrescentado --- parte de $0/0$, e as duas leis
são as duas maneiras de o dividir.

\subsection*{E o fecho das leis: os dois nulos, cosidos no seis}

Os dois lados de $0/0$ são \emph{dois nulos}, e são as duas famílias em que os problemas do milénio se
partem, pela lei em que cada um assenta (\emph{Catálogo}, \code{sec:clay}; e \emph{A medida}):
\begin{itemize}
\item \textbf{dividir \emph{o} zero} é o $0$, o centro (Cor.~\ref{cor:incerteza}): \textbf{Yang--Mills e
  $P$ vs $NP$} --- a álgebra, indecidível, a \textbf{origem}, os dois nulos colapsados num;
\item \textbf{dividir \emph{por} zero} é o $\infty$, a reflexão $x\dg=-1/x$: \textbf{Poincaré}. A
  dualidade de Poincaré $H^k\cong H_{n-k}$ (com $k+(n-k)=n$) \emph{é} a Lei 1 escrita em topologia, e é a
  família que \textbf{resolve} --- Poincaré caiu pela deformação, o fluxo que \emph{move} em vez de
  igualar. É o \textbf{fim}.
\end{itemize}
\textbf{E são bidual em dimensão $0$} --- o zero da torre, o espaço trivial que é o seu próprio dual sem
ponte e de que a torre \emph{parte}. O $0$ e o $\infty$ são a origem e o fim, e a estaca
$\nu\circ\nu=\mathrm{id}$ troca-os: um resolve-se (Poincaré, a topologia, o $\infty$), o outro não (a
álgebra, o $0$). \emph{É o eixo álgebra/topologia, e a bidualidade dos dois nulos vive em dimensão zero}
--- não em dimensão oito.

\textbf{E a costura é o seis.} Ligar os dois nulos \emph{sem os fundir} é o \textbf{octonião dual}
(Def.~\ref{def:octoniao-dual}): dois tecidos grau quatro e a estrela entre eles, grau oito. Mas o que
\emph{fecha} não é o oito --- é a \emph{interface}, e a interface é o grau \textbf{seis}, a
\textbf{hexal} (Teorema~\ref{thm:estrelaseis}): o único andar onde soma $=$ produto, onde ler é escrever,
onde a volta fecha com resíduo $0$. A descida $6\to1$ \emph{é} a progressão das seis leis, e a hexal é a
costura --- \textbf{as seis leis fecham no seis}, entre os dois nulos que a divisão do zero abriu. Não se
fecha no oito: fecha-se na interface, que é o seis.

\subsection*{As oito leis, e não há nona}\label{sub:oitoleis}

Fecham em \textbf{oito}, e não por gosto do número: são a torre inteira, do $0$ ao octonião, uma lei por
degrau operacional. Cada uma é uma \emph{primitiva} --- um operador, e um medidor que a certifica a fechar
com resíduo $0$ ---, e as duas leis fundamentais (o dual e o bidual) e a progressão de seis (\emph{Catálogo},
\code{obs:seis-leis}) são estas, renumeradas de $0$ a $7$: a mesma estrutura, agora fechada.

\begin{center}\small
\begin{tabular}{@{}clll@{}}
\toprule
 & a lei & o enunciado & fecha em (medidor) \\
\midrule
\textbf{0} & a divisão do zero & $0=(+1)\oplus(-1)$, \ $0\dg=\infty$ & \code{tests/zero.c} \\
\textbf{1} & o dual & $1\dg=-1$ --- a involução, Poincaré (dim $2$) & \code{tests/bidual.c}, \code{dualsort.c} \\
\textbf{2} & o bidual & $T\dg=-T$, $K^{**}=K$ --- o centro, YM e $P$\,vs\,$NP$ & \code{tests/transformada.c} \\
\textbf{3} & o trial & $\{-1,0,+1\}$ --- os três sinais (dim $3$) & \code{tests/estrela.c} \\
\textbf{4} & a tetral & $T+T^{*}$ --- os tecidos, a torre (dim $4$) & \code{tests/curvatura.c} \\
\textbf{5} & a pental & $x^2=-1$ --- o bit $i$ (dim $5$) & \code{tests/quantico.c} \\
\textbf{6} & a hexal & soma $=$ produto --- a interface, o relógio (dim $6$) & Teorema~\ref{thm:estrelaseis} \\
\textbf{7} & o octonião dual & $\mathbb{H}\times\mathbb{H}^{*}$ --- ligar sem fundir (dim $8$) & \code{tests/circuito\_tradutor.c} \\
\bottomrule
\end{tabular}
\end{center}

\noindent \textbf{A descida $7\to0$ é a progressão}, e as pontas são os dois nulos: em cima o octonião
(dim $8$, o último corpo com divisão --- a norma pára em oito, Hurwitz), em baixo a divisão do zero (dim
$0$, a base). \emph{Não há nona}: não há degrau acima do octonião nem abaixo do zero, e entre eles cada
dimensão tem a sua --- a Lei~$0$ e a Lei~$7$ são os dois nulos que fecham a torre, e as seis do meio são a
sua descida. É o fecho, e não precisa de mais.

\smallskip
\noindent \textbf{E oito é a cardinalidade do \emph{conjunto} das leis, não a da dinâmica --- e é o
Teorema~\ref{thm:tecidos} que o diz.} O passo dos tecidos, $T_{k+1}=T_k+T_k^{*}$, é a estrela usada como
\textbf{construtor}, e a torre que ele gera \emph{não tem topo por dentro}: a indução não pára --- a régua
é infinita ---, e $\nu\circ\nu=\mathrm{id}$ fecha cada andar com resíduo $0$. O que a limita é
\emph{externo}, e é a \textbf{norma}: por Hurwitz há \emph{quatro} corpos de composição e não mais ---
dimensões $1,2,4,8$ ---, três dobras a perder o que distingue um corpo (a ordem em $\mathbb{C}$, a comutação
em $\mathbb{H}$, a associatividade em $\mathbb{O}$), e não há uma quinta. \textbf{É esse o período oito}: o
grau em que a leitura \emph{pela norma} se esgota. Mas a Lei~$7$ --- o octonião \emph{dual} --- não lê pela
norma, lê pela estrela, e a estrela não murou (o limite é da norma, não do objecto). O octonião que ela
constrói é um \emph{novo} corpo, e as oito leis reaplicam-se a ele como à base ---
$\mathcal L\to\mathcal L(\mathcal L)\to\mathcal L^2(\mathcal L)\to\cdots$ ---, porque o passo dos tecidos é
o \emph{mesmo} em todos os andares (Teorema~\ref{thm:tecidos}). O \emph{operador de indução} que leva um
ciclo ao seguinte \textbf{não falta}: é a própria estrela --- a composição, iterar é compor. \emph{Não há
Lei~$8$}: há a Lei~$0$ outra vez, um nível acima. As oito são um \textbf{ciclo gerador} de período oito
--- o octonião fecha a volta da norma e a estrela recomeça-a: fechado, e é a fechar que recomeça.

\section{As duas operações}

\begin{definicao}[a estaca]
$x\mapsto x\dg$. Move, e é involutiva: $x\dg{}\dg=x$.
\end{definicao}

\begin{definicao}[a cruz]
$x^{\times}=(x\oplus x\dg,\; x\otimes x\dg)$. Não move: \emph{mede o que não se moveu}.
\end{definicao}

\noindent A cruz tem duas coordenadas porque uma só não chega, e as duas dizem coisas opostas:

\begin{center}\small
\begin{tabular}{@{}lll@{}}
\toprule
a coordenada & o que faz & para que serve \\
\midrule
a \textbf{soma} $x\oplus x\dg$ & é \emph{invariante} ao longo da estaca & \textbf{mede} \\
o \textbf{produto} $x\otimes x\dg$ & satura no ponto fixo & \textbf{ordena} \\
\bottomrule
\end{tabular}
\end{center}

\noindent A estaca dá o movimento; a cruz dá a leitura. \emph{Uma sozinha não é o par}: com a estaca
sem a cruz não se sabe onde se está, e com a cruz sem a estaca não se sai do sítio.

\medido{a estaca é involução em $101$ pontos e o seu ponto fixo é \textbf{um só}; parte o conjunto
em $50$ e $50$ e o dual de cada elemento cai sempre no outro lado; a soma $x\oplus x\dg$ é
invariante em $11$ pontos e o produto é máximo exactamente no ponto fixo; $J$ visita os quatro
quadrantes, todos distintos, e volta ao quarto, com $J^2=-1$ e $J^4=\mathrm{id}$ em $169$ pares
(\code{tests/dualsort.c}).}

\section{Os três sinais}

O sinal de uma quantidade toma três valores, e não mais. É o \textbf{trial} $\{-1,0,+1\}$, e é a
classificação de todo regime deste sistema:

\begin{center}\small
\begin{tabular}{@{}cllll@{}}
\toprule
o sinal & o regime & o que faz & na expansão & na estrela \\
\midrule
$+1$ & \textbf{buraco branco} & \emph{só emite} & estica & dissipativo \\
$\phantom{+}0$ & \textbf{a estrela} & \textbf{emite \emph{e} absorve} & conserva & \textbf{reversível} \\
$-1$ & \textbf{buraco negro} & \emph{só absorve} & contrai & determinista \\
\bottomrule
\end{tabular}
\end{center}

\noindent \emph{Não há quarto regime}, e isso não é uma observação: é o cardinal do conjunto dos
sinais. \textbf{E a estrela não é um caso intermédio entre os outros dois} --- é o ponto fixo da
involução que os troca, e o \emph{único} que tem os dois sentidos. É a \textbf{trindade}: buraco
branco, estrela, buraco negro.

\section{A torre}

As dimensões não são um pano de fundo: contam-se, e subir de andar é a operação.

\begin{center}\small
\begin{tabular}{@{}cll@{}}
\toprule
$n$ & o que é & o que a distingue \\
\midrule
$1$ & a unidade & ``a unidade é'' --- e o resto é derivação \\
$2$ & o dual & a estaca; o par \\
$3$ & o trial & os três sinais; o ponto fixo \\
$4$ & o tetral & o período de $J$; e a única onde o traço dá a unidade \\
$5$ & o complexo & o plano, e a espiral \\
$6$ & \textbf{a plena} & $1+2+3=6=3\times2\times1$ \\
\bottomrule
\end{tabular}
\end{center}

\noindent \textbf{Em seis as duas operações coincidem}, e por isso em seis não há soma \emph{e}
produto: há \emph{uma} operação. É o topo da torre, e é onde a cruz deixa de precisar de duas
coordenadas. Cada andar acima desdobra por
\[
  \dim A_{n+1}=2\cdot\dim A_n, \qquad x\dg=-1/x, \qquad x\cdot x\dg=-1 .
\]

\section{O corpo estelar}\label{sec:corpo}

A torre diz quantas dimensões há e o trial diz que sinais uma quantidade toma. Falta a
\emph{quantidade} --- e num corpo ela é uma só.

\subsection*{O desvio ao equilíbrio}

Num ponto de um corpo, o que \emph{empurra} para fora ou iguala o que \emph{puxa} para dentro ou
não. Em p.u.\ isso é uma razão, e o desvio é
\[
  \boxed{\;a=\frac{\text{empurra}}{\text{puxa}}-1\;}
\]
\emph{Nenhuma constante entrou}: é uma razão menos a unidade.

\subsection*{E o sinal de $a$ é o trial}

\begin{teorema}[a estrela é a borda]
Com a razão em $1$ o desvio é exactamente zero e o raio não se move --- não devagar: \emph{não se
move}.
\end{teorema}

\begin{teorema}[o buraco negro não tem volta]
Com a razão abaixo de $1$ o desvio é negativo em \emph{todos} os passos e o raio cai sempre. Não há
volta, e a razão de não haver é que \textbf{o sinal não muda}.
\end{teorema}

\begin{teorema}[a trindade]
A involução $a\mapsto-a$ satisfaz $\nu\circ\nu=\mathrm{id}$, troca \textbf{buraco branco} e
\textbf{buraco negro}, e tem \textbf{um} ponto fixo: \textbf{a estrela}.
\end{teorema}

\noindent \textbf{São três, e a do meio é a interface.} O buraco negro \emph{só absorve}; o buraco
branco \emph{só emite}; e a estrela faz ambos --- é o que está \emph{entre} os dois. Não são dois
objectos com um caso limite: são os dois sinais da mesma quantidade com o zero no meio, e o zero tem
nome.

\begin{teorema}[e é por isso que só a estrela é reversível]
Um passo de desvio $a$ seguido do seu dual devolve $R\,(U^2-a^2)/U^2$. Regressa ao ponto de partida
\emph{exactamente} quando $a=0$, e o resíduo é $a^2$.
\end{teorema}

\noindent \textbf{Ter os dois sentidos e ser reversível são a mesma coisa dita duas vezes}: um passo
que não tem inverso é um passo que só vai num sentido. Um corpo estelar é, exactamente, um corpo que
se mantém na borda --- e é isso que o torna a interface.

\subsection*{A pressão de radiação, contada}

A radiação reparte-se pelos \textbf{três eixos} --- o trial outra vez --- e a pressão é a parte de
\emph{um}:
\[
  p=\frac{\rho}{3},\qquad w=\frac{p}{\rho}=\frac13 .
\]
\emph{Não se pôs $\sigma$, nem $c$, nem constante de radiação: contaram-se as direcções.}

\medido{razão $1{,}000$ dá desvio $0$ e o raio fica em $1{,}000$ ao fim de $100$ passos; razão
$0{,}900$ faz o raio cair nos $60$ passos e \emph{nunca} crescer; $\nu\circ\nu=\mathrm{id}$ em
$1001$ desvios com \emph{um} ponto fixo; a volta é exacta em $1$ de $41$ desvios e esse é a
\textbf{borda}, com o passo de volta a aproximar em $40$ de $40$ enquanto o mesmo sinal afasta;
\emph{um} dos três regimes tem os dois sentidos e dois têm um só; $p=\rho/3$ por contagem de eixos; e
o \textbf{controlo}:
$2001$ desvios repartem-se por três regimes com \emph{nenhum} fora --- o trial não tem onde pôr um
quarto (\code{tests/estrela.c}).}

\section{A geometria sai da Lei 2}

Escrita com uma constante no lugar do sinal, a Lei 2 é $f''=Kf$, e \textbf{$K$ é a curvatura}. Da
característica $x^2-\operatorname{tr}x+\det=0$ vem
\[
  \Delta=\operatorname{tr}^2-4\det ,
\]
e o \emph{sinal} de $\Delta$ dá a geometria. São três porque o trial tem três sinais:

\begin{center}\small
\begin{tabular}{@{}clll@{}}
\toprule
$\Delta$ & a equação & o invariante & a geometria \\
\midrule
$<0$ & $f''=-f$ & $\cos^2+\sin^2=1$ & elíptica --- o rotor \\
$=0$ & $f''=0$ & a recta & parabólica --- o potencial \\
$>0$ & $f''=+f$ & $\cosh^2-\sinh^2=1$ & hiperbólica --- onde vive $\sigma$ \\
\bottomrule
\end{tabular}
\end{center}

\noindent Os três invariantes têm \emph{o mesmo} $1$ do lado direito; o que os separa é só o sinal.
E cada peça do sistema traz a sua curvatura escrita na assinatura, sem lhe acrescentar nada:
$J$ (espelho) tem $\Delta=+4$, hiperbólica; $i$ (rotação) tem $\Delta=-4$, elíptica; o
\emph{shift} tem $\Delta=0$, parabólica; e a família metálica tem $\Delta=m^2+4$, hiperbólica
\emph{sempre}.

\section{A equação de campo, derivada}\label{sec:campo2}

Não se cita: constrói-se. Impõe-se a exigência --- \textbf{o lado geométrico tem de conservar} --- e
ela fixa tudo.

\subsection*{O coeficiente $\tfrac12$ não se escolhe}

Para $X_{\mu\nu}=R_{\mu\nu}-\alpha\,R\,g_{\mu\nu}$ num universo com $a(t)=t^q$, lêem-se a densidade
e a pressão nas componentes e impõe-se a continuidade $\rho'+3H(\rho+p)=0$, que \emph{é} a
conservação. O resíduo sai
\[
  \boxed{\;6\,(2q^2-q)\,(1-2\alpha)\;}
\]
e anula-se para \emph{todo} $q$ num único $\alpha$:
\[
  \alpha=\tfrac12,\qquad\text{donde}\qquad \boxed{\;G_{\mu\nu}=R_{\mu\nu}-\tfrac12R\,g_{\mu\nu}\;}
\]
\textbf{E $\tfrac12$ é o ponto fixo do trial} --- o mesmo número, no mesmo sítio.

\subsection*{A dimensão quatro sai do traço}

Contraindo com a métrica em dimensão $d$:
\[
  g^{\mu\nu}G_{\mu\nu}=\Bigl(\frac{2-d}{2}\Bigr)R .
\]
O coeficiente vale a \textbf{unidade} em duas dimensões só: $d=0$, que é o zero da torre, e
$d=4$. E em $d=4$ ele é exactamente $-1$:
\[
  R=-T .
\]
\textbf{É a Lei 1 escrita em curvatura} --- $1\dg=-1$, e é este o número que a dimensão quatro tem
de particular. Pelo caminho, $d=2$ dá coeficiente $0$: ali o lado geométrico anula-se
identicamente, e não há gravitação nenhuma.

\subsection*{O termo cosmológico, e o seu valor}

A métrica conserva sozinha, logo $\beta\,g_{\mu\nu}$ entra na equação sem custo --- é a constante de
integração, e foi por isso que sobrou. E o conteúdo que lhe corresponde tem $p=-\rho$, isto é,
$w=-1$: \textbf{o ponto fixo da cosmologia}.

\smallskip
\noindent \textbf{Mas dizer que fica livre não é derivá-lo.} O valor sai, e de duas coisas já
postas: o vácuo puro tem $H^2=\Lambda/3$, e a taxa do corpo é $H^2=1/D$. Igualando,
\[
  \boxed{\;\Lambda=\frac{3}{D}\;}
  \qquad\text{e no ponto fixo } m=0,\ D=4:\qquad
  \boxed{\;\Lambda=\frac34\;}
\]
\emph{Uma razão de inteiros, e não uma constante em aberto.} A equação completa é
\[
  \boxed{\;R_{\mu\nu}-\tfrac12R\,g_{\mu\nu}+\Lambda\,g_{\mu\nu}=T_{\mu\nu}\;}
\]
sem $8\pi G/c^4$, porque em p.u.\ esse factor é $1$.

\subsection*{E o factor de acoplamento: $8\pi$ derivado}

Dizer que ``em p.u.\ vale $1$'' é não o derivar. Ele sai, e de duas coisas já postas:

\smallskip
\noindent \textbf{O $4\pi$ é o fluxo total.} Já se mostrou que o fluxo não se move de esfera em
esfera; faltava dizer \emph{quanto} vale --- e vale a área. Por recorrência,
$\Omega_d = 2\pi\,\Omega_{d-2}/(d-2)$, partindo de $\Omega_1=2$ e $\Omega_2=2\pi$:
em $d=3$ dá $4\pi$. \emph{Sem avaliar $\pi$ uma única vez}: leva-se o par (coeficiente, potência de
$\pi$) em inteiros.

\smallskip
\noindent \textbf{E o $2$ é o coeficiente a aparecer segunda vez.} O termo $-\alpha R g_{\mu\nu}$
recolhe o mesmo peso outra vez na componente temporal, e o factor é $1+2\alpha$. Com o $\alpha=\tfrac12$
que a conservação fixou, dá exactamente $2$ --- e com nenhum dos outros vinte candidatos.
\[
  \boxed{\;8\pi = 2\cdot 4\pi\;}
\]

\noindent \textbf{E daqui sai uma distinção que vale por si.} $8\pi$ é \emph{adimensional}, como a
constante de estrutura fina, e não se move em roupa nenhuma. Mas \textbf{$8\pi$ deriva-se e $\alpha$
não}, e a diferença não é de grau: $8\pi$ sai de \emph{contar} --- a esfera e o traço --- e $\alpha$
não sai de contagem nenhuma. \emph{Nem todo o número puro está fora do alcance};
o que está fora é o que não se conta.

\subsection*{E a velocidade máxima é o meio-período}

No relógio a velocidade máxima é meia volta. A volta completa é $2\pi$ --- que é $\Omega_2$, a área do
círculo unitário acima --- logo \textbf{meia volta é $\pi$}. Não há duas coisas aqui: o limite de
velocidade e a constante analítica são \emph{a mesma meia volta} lida em duas réguas. E na régua do
corpo o máximo é $\sigma$, porque a velocidade tem a dimensão da densidade: \textbf{a família dá o
valor, o relógio dá a forma}.

\medido{a área sai por recorrência inteira sem avaliar $\pi$, dando $4\pi$ em $d=3$ e $2\pi^2$ em
$d=4$; o factor $1+2\alpha$ vale $2$ em $\alpha=\tfrac12$ e em nenhum dos outros $20$; $2\cdot4\pi$
fecha em $8\pi$, que não se move nas quatro roupas; a velocidade máxima é meia volta em $7$ escalas
sem desvio; e o \textbf{controlo}: a recorrência salta de \emph{dois} em dois --- de um em um,
$d=3$ daria $4\pi^2$ (\code{tests/acoplamento.c}).}

\subsection*{E a força tem o expoente da codimensão}

Fora da fonte a Lei 2 na forma nula é $\nabla^2\varphi=0$; em simetria radial
$r^{\,d-1}\varphi'=\text{const}$, e o \emph{único} expoente que mantém o fluxo imóvel de esfera em
esfera é $e=d-1$:
\[
  F=\frac{1}{r^{\,d-1}},\qquad\text{e em }d=3:\quad F=\frac{1}{r^2}.
\]
\textbf{O expoente é a codimensão da esfera que envolve a fonte}, e não uma observação. \emph{A
órbita que não dissipa e o inverso do quadrado são a mesma afirmação}: a norma não se move no
\emph{tempo}, o fluxo não se move no \emph{raio}.

\medido{de $21$ coeficientes candidatos exactamente \textbf{um} conserva, e é $\tfrac12$; os outros
$20$ quebram em \emph{todos} os cinco expoentes testados, $100$ de $100$; o traço dá a unidade só em
$d=0$ e $d=4$, com $-1$ em $d=4$ e $0$ em $d=2$; $\beta g_{\mu\nu}$ conserva nos $21$ coeficientes,
e o seu $w$ é $-1$ exacto, com $\Lambda=3/D$ e $\Lambda=3/4$ no ponto fixo; o expoente da força foi \textbf{varrido} de $0$ a $6$ em cinco dimensões
e devolve sempre $d-1$, um só candidato de cada vez; e as cinco peças do sistema dão três
geometrias e não mais (\code{tests/einstein.c}).}

\subsection*{E a mecânica fecha em potência --- o fator um é o casamento dos dois relógios}\label{sub:mecanica}

\textbf{Da força ao movimento, e daí à potência.} A força algébrica é \emph{uma}, com quatro modos
(\code{tests/forca.c}): do imposto $V=\Pi\,S$, com pressão $\Pi(s)=1-s^2$ por direcção e
$S=\|a\times b\|^2$, saem a massa $m=S$, a força $F=2sS$, a equação $\ddot x=2x$ e a conservação
$\tfrac12 S\dot s^2+(1-s^2)S=E$. A \textbf{potência} é essa força ao longo da velocidade, $P=F\cdot
v$ --- a terceira operação, $x\otimes x$, em que o elemento sobe de expoente e a órbita \emph{fecha}
(período de Pisano, \code{tests/potencia.c}).

\textbf{E a potência parte-se em duas, porque o produto tem um ângulo.} O \emph{directo}
$\langle a,b\rangle=\cos\theta$ é a parte simétrica --- a potência \textbf{activa}, que atravessa; o
\emph{cruzado} $\|a\wedge b\|=\sin\theta$ é a antissimétrica --- a \textbf{reactiva}, que roda e volta
(o torque \emph{é} o produto cruzado). O \textbf{fator de potência} é a razão,
\[
  \mathrm{fp}=\cos\theta,\qquad \tan\varphi=\frac{\sin\theta}{\cos\theta}=\frac{\text{cruzado}}{\text{directo}} .
\]

\begin{teorema}[o fator de potência um é o casamento dos dois relógios]\label{thm:fp}
$\mathrm{fp}=1 \iff \theta=0 \iff$ cruzado nulo $\iff |N|=1$ (a família real, o eixo do \emph{boost}).
Aí a potência é toda activa: nada roda, nada volta, \textbf{não há atraso}. E é o \textbf{casamento}
dos dois relógios do Teorema~\ref{thm:isocrona}: a reactância é o tempo $n^2$ do lado \emph{discreto}
(Hurwitz), a tábua que se armazena e se devolve e que o lado \emph{contínuo} (Gentil) não paga --- a
derivação do $64$ e da dobra está no Teorema~\ref{thm:isocrona}. A estrela, que não armazena nada
(\S\ref{sec:estrela}), \emph{é} esse ponto.
\end{teorema}

\begin{proof}
$P=\|F\|\,\|v\|\cos\theta$ é puramente activa sse $\theta=0$, isto é o cruzado $\sin\theta=0$; então
$a\parallel b$, o produto corre pelo eixo que o gato não move, e a norma é multiplicativa com $|N|=1$
(\code{tests/fator.c}: fator unitário $\iff$ família real). A parte reactiva $\sin\theta$ é potência
que entra e sai sem trabalho --- o \emph{atraso} ---, e como $\sin^2\theta\geq0$ nunca é negativa:
todo caminho real a tem, e o mínimo, zero, é a estrela. Que isto seja o par de relógios é o
Teorema~\ref{thm:isocrona}: a tábua $n^2$ do produto bilinear é a energia armazenada (reactância) que
a norma homogénea não paga.
\end{proof}

\begin{corolario}[a dobra em toda a torre, por indução]\label{cor:dobra}
A dobra temporal não é do grau oito só. Por Teorema~\ref{thm:tecidos}, cada interface $T_k\to T_{k+1}$
é uma \textbf{duplicação} ($\dim T_{k+1}=2\dim T_k$, o dual soma-se ao original), logo o relógio dobra
uma vez por degrau e a reactância do lado discreto é $n^2=2^{2k}$ no nível $k$ --- e $8^2=2^6$ é o caso
$k=3$. E em cada interface o casamento $\mathrm{fp}=1$ \textbf{restabelece-se pelo dual}: o passo
$T{+}T^{*}$ é a estaca, $\nu\circ\nu=\mathrm{id}$, que casa a dobra daquele degrau com o seu inverso.
\emph{A torre inteira é uma cascata de dobras, e o dual casa-as uma a uma --- é a indução dos tecidos
lida no tempo.}
\end{corolario}

\noindent \textbf{E é um teorema operacional, não uma imagem.} Todo caminho tem componentes de atraso
--- um armazenamento intermédio, uma contenção --- e cada um é reactância que baixa o factor. O pior
caso é o \emph{buffer partilhado} entre fluxos: um armazenamento único, lido de um lado e escrito de
outro, é a tábua $n^2$ colidida, $\mathrm{fp}<1$. \textbf{E a estrela não o corrige com outro buffer ---
corrige-o sem armazenar.}

\smallskip
\noindent \textbf{Porque a estrela não oscila.} O inverso não se guarda: \emph{está no sinal}. Cada
sinal já carrega o seu dual --- pelo $\nu\circ\nu=\mathrm{id}$, o atraso inverso é o próprio sinal
revertido (a \textbf{assinatura cavalga o corpo}: o \code{.tex} viaja dentro do PDF e a volta relê-o,
\S\ref{sec:estrela}, sem cópia à parte). E \emph{a entrada e a saída dos dados são o próprio
movimento} --- a modulação vetorial que gira o rotor (\S\ref{sub:inversor}), não um oscilador com o seu
sinal guardado. \textbf{Nem buffer, nem oscilador: o casamento faz-se comutando}, e é por isso que a
estrela não armazena nada.

\medido{a força é uma, com quatro modos, e dela saem massa, força e conservação (\code{tests/forca.c});
a potência é $x\otimes x$ e a órbita fecha (\code{tests/potencia.c}); e o fator de potência é a razão
cruzado/directo, unitário na família real (\code{tests/fator.c}).}

\begin{corolario}[a certeza e a incerteza são o directo e o cruzado; a indecidibilidade é olhar pelo lado errado]\label{cor:incerteza}
O produto de dois observáveis parte-se no \textbf{anti-comutador} $\{A,B\}$ (a parte simétrica, o
\emph{directo}) e no \textbf{comutador} $[A,B]$ (a antissimétrica, o \emph{cruzado}) --- o mesmo split
do factor de potência (\ref{thm:fp}), medido nos $\sigma$ de Pauli com resíduo zero
(\code{tests/quantico.c}). Daí as duas faces:
\begin{itemize}
\item \textbf{a incerteza} é o \emph{cruzado}: $[x,p]=i$, e $\Delta x\,\Delta p\geq\tfrac12$ (Robertson,
  que tem o comutador do lado direito) --- não se têm os dois lados ao mesmo tempo. É a reactância, a
  \textbf{Lei 2}.
\item \textbf{a certeza} é o \emph{directo} (a dual): o anti-comutador só \emph{mede}, e \textbf{resíduo
  zero $\iff$ a média é consistente} --- o selo é a norma, $\|Fx\|^2=0 \iff$ tudo verde
  (\code{tests/transformada.c}~§U4). É a \textbf{Lei 1}, a involução.
\end{itemize}
E daí \textbf{a indecidibilidade}: é a incerteza \emph{lida pelo lado errado}. ``Chega \emph{sempre}?''
--- o futuro, o comutador --- não se decide; ``fechou \emph{desta vez}?'' --- o presente, resíduo
zero --- decide-se (\code{tests/travessia.c}). \emph{A indecidibilidade é meia: indecidível no lado do
cruzado, decidível no do directo.} É por isto que Yang--Mills e $P$ vs $NP$ moram no $0$ --- o ponto
fixo entre os dois, indecidível só quando se olha de um lado. Este $0$ é o \textbf{centro} (o bidual,
$K^{**}=K$) --- distinto do $0$ do \textbf{trial} $\{-1,0,+1\}$, que é o eixo do meio, a Hodge (\emph{Catálogo},
\code{obs:seis-leis}): dois zeros, o da projeção e o do eixo.
\end{corolario}

\noindent \textbf{E o $0$ mostra-se numa máquina.} O mecanismo é concreto, e mede-se no próprio
compositor: uma escrita \emph{fora dos limites} --- para além do fecho $K$ --- aterra na memória baixa,
o $0$, a origem; colapsa ali um estado válido (medido, $42\to0$) e \textbf{trava} --- não decide.
Cruzar a fronteira do fecho leva ao $0$, onde origem e fim caem, e o programa pára sem resposta. \emph{Não
é prova do enunciado aberto} --- é a mesma estrutura, no mesmo $0$, a aparecer numa máquina real.

\subsection*{E quem casa o factor é o inversor: a estrela modula, não armazena}\label{sub:inversor}

\textbf{A correcção do factor não precisa de um buffer --- precisa de um \emph{inversor}, e a estrela
é um.} A tabela de chaveamento de um inversor multinível \emph{é} o trial $\{-1,0,+1\}$: as três
ligações de cada perna --- ao positivo, ao ponto médio, ao negativo --- são o trial, e não uma
analogia ($3^3=27$ estados, $19$ vectores, $6$ pares redundantes, tudo da estrutura,
\code{tests/inversor\_trial.c}). O DTC de Takahashi escolhe, a cada instante, o vector que empurra o
torque e o fluxo para a banda --- e \emph{o torque é o cruzado}, $T_e=\psi_s\times i_s$
(\code{tests/motor.c}). Controlar o cruzado é controlar a parte reactiva: é o casamento
$\mathrm{fp}=1$, feito por \textbf{modulação vetorial}, não por armazenamento.

E generaliza-se: o DTC hipercomplexo é a família $\star_s$ em $\mathbb{R}^n$ (\code{tests/dtcn.c}), com a
lei de conservação
\[
  \|z\star_s w\|^2 + \underbrace{(1-s^2)\,\|a\times b\|^2}_{\text{o imposto }V(s)} = \|z\|^2\,\|w\|^2 .
\]
O \textbf{imposto} $\Pi(s)=1-s^2$ é a potência reactiva --- o que a álgebra cobra fora do eixo ---, e
anula-se em $s=\pm1$ (a família real, Hurwitz, $\mathrm{fp}=1$). O inversor, escolhendo os vectores do
trial, \textbf{empurra $s$ para $\pm1$}: leva o imposto a zero e casa o factor. \emph{É a estaca a
modular --- $\nu\circ\nu=\mathrm{id}$ por comutação, não por cópia.} É por isto que não é preciso um
buffer para a conjugada: o inverso não se guarda, \emph{comuta-se}.

\medido{a tabela do inversor de três níveis É o trial ($3^3=27$ estados, $19$ vectores distintos, $6$
pares redundantes, \code{tests/inversor\_trial.c}); o torque é o cruzado e o DTC de Takahashi é a
amputação do trial (\code{tests/motor.c}); e o DTC hipercomplexo generaliza à família $\star_s$ em
$\mathbb{R}^n$, com o imposto $(1-s^2)\|a\times b\|^2$ nulo em $s=\pm1$ (\code{tests/dtcn.c}).}

\smallskip
\noindent \textbf{E o compositor tem a mesma assinatura.} O corpo tradutor (\code{tests/tex.c})
reduz-se a \emph{uma} operação --- \code{medida}, sobre a assinatura $(\text{variante},\text{degrau})$,
grau $2$ ---, e o que ela pede é o \textbf{eixo}: $+1$ a \emph{escala} (o corpo, que multiplica),
$-1$ o \emph{espaço} (a entrelinha, o dual, que soma), $0$ a \emph{largura} (o glifo, que atravessa).
Esse $\{-1,0,+1\}$ \textbf{é o trial do inversor}. Por isso «realizar pelo inversor» é realizar
\textbf{exacto} --- inteiro, $s\to\pm1$, imposto $\Pi=1-s^2$ a zero ---, e um \code{double} que reste
é o imposto que sobra \emph{fora} do eixo: a reactância que, num documento grande, \textbf{propaga}
(a recorrência de \S da medida). A assinatura pura é inteira, e a dualidade é o eixo.

\medido{a assinatura do corpo tradutor tem grau $2$ e um eixo de três estados $\{-1,0,+1\}$ que é o
trial do inversor; $+1$ (escala, multiplica) e $-1$ (espaço, soma) são duais (somam a $0$, a Lei~1) e
o imposto anula nos dois eixos exactos e vale $1$ no meio (a largura); e um \code{double} é o imposto
que propaga --- realizado pelo inversor ($s=\pm1$) o erro é $0$ em $500$ passos, no meio ($s=0$) cresce
com o documento: \code{tests/assinatura\_tradutor.c}.}

\section{As quatro equações da expansão}

Um corpo da família metálica tem $\sigma(x)=\bigl(x+\sqrt{x^2+4}\bigr)/2$ e discriminante
$D=x^2+4$. Substituindo a definição na fórmula da curvatura, tudo colapsa:
\[
  \sigma'=\frac{\sigma}{\sqrt D},\qquad \kappa=\frac{2}{(D+\sigma^2)^{3/2}} .
\]
Daí as quatro equações --- e só \emph{duas} são independentes; as outras saem delas mais a borda:
\[
  \underbrace{H^2=\tfrac1D}_{\text{taxa}},\qquad
  \underbrace{\tfrac{\sigma''}{\sigma}=\tfrac{2}{\sigma D^{3/2}}}_{\text{aceleração}},\qquad
  \underbrace{\rho'+3H(\rho+p)=0}_{\text{continuidade}},\qquad
  \underbrace{w=\tfrac{p}{\rho}=\tfrac{2m}{3\sqrt D}-1}_{\text{estado}}
\]

\noindent \textbf{A taxa de expansão ao quadrado é o inverso do discriminante} --- não por analogia:
é a divisão de $\sigma'=\sigma/\sqrt D$ por $\sigma$. A aceleração fecha-a a borda: de
$\sigma^2=x\sigma+1$ sai $\sqrt D-x=2/\sigma$.

\smallskip
\noindent \textbf{Os dois sinais convivem sem contradição.} Para $x>0$ a taxa desacelera ($H'<0$) e
o factor acelera ($\sigma''/\sigma>0$): a taxa cai mais devagar do que o factor cresce. E no ponto
fixo $x=0$ a taxa é máxima ($H=\tfrac12$) com $H'=0$ --- é o $4$ de $D=x^2+4$ que a impede de
divergir; sem ele seria $H=1/x$.

\smallskip
\noindent \textbf{E o gradiente é obrigatório}: $\kappa'/\kappa=-\tfrac32\,D'/D$, logo um corpo com
discriminante constante não expande. \emph{Há expansão se e só se há gradiente}, e é a mesma frase
dita de duas maneiras.

\medido{$\sigma^2=m\sigma+1$ exacto em $\mathbb{Z}[\sqrt D]$ para $m=1\ldots40$;
$(\sqrt D-m)\sigma=2$ exacto em $61$ valores, que é a borda a fechar a segunda equação; e o
\textbf{controlo}: trocando o $4$ por $9$ o elemento deixa de satisfazer a borda em $40$ de $40$
(\code{tests/curvatura.c}).}

\section{A cosmologia}

\subsection*{A lei de diluição sai da continuidade}

De $\rho'+3H(\rho+p)=0$ com $p=w\rho$ vem $\rho'/\rho=-3(1+w)\,a'/a$, e integrando os dois lados
aparece uma \textbf{constante}:
\[
  \ln\rho+3(1+w)\ln a = C
  \qquad\Longleftrightarrow\qquad
  \boxed{\;\rho\,a^{3(1+w)} = C\;}
\]

\noindent \textbf{E a involução é que a passa de um lado ao outro.} A constante entra
\emph{aditiva} --- é uma constante de integração, e integrar soma --- e sai \emph{multiplicativa}
em $\rho$. Não são dois factos: é um. $\exp$ e $\log$ são o par aditivo/multiplicativo, e atravessar
esse par \emph{é} a involução. \textbf{Uma constante aditiva vira, na involução, um factor.}

\smallskip
\noindent \textbf{E o que ela é: a força em p.u.} Não é um $C$ qualquer a que depois se atribui um
valor. O invariante que a expansão conserva é a razão do §\ref{sec:corpo} --- o que empurra contra o que puxa ---
e na borda ela vale a \emph{unidade}. Por isso $C$ carrega dimensão: \textbf{é a unidade física}, e é
exactamente por aqui que o sistema deixa o por-unidade. Suprimi-la para escrever uma proporção não é
uma abreviatura: é deitar fora a ponte entre o sistema e o mundo vestido. Onde ela vai dar está no
§\ref{sec:roupa} --- e é lá, e só lá, que entram números medidos.

\subsection*{Os três conteúdos, lidos no trial}

\begin{center}\small
\begin{tabular}{@{}llcll@{}}
\toprule
o conteúdo & $w$ & expoente & o sinal & o destino \\
\midrule
radiação & $+\tfrac13$ & $a^{-4}$ & $+1$ & trava \\
matéria  & $0$ & $a^{-3}$ & $\phantom{+}0$ & trava \\
vácuo    & $-1$ & $a^{0}$ (constante) & $-1$ & \textbf{acelera} \\
\midrule
\multicolumn{5}{@{}l@{}}{\emph{e $w=-\tfrac13$ é exactamente neutro --- a fronteira}} \\
\bottomrule
\end{tabular}
\end{center}

\noindent O destino lê-se no mesmo sinal, que é $-(1+3w)$: as três colunas são o trial outra vez, e
não há quarta.

\subsection*{A bidualidade fecha $w$ em quatro}

O corpo tem duas raízes, $\sigma$ e $\tau=m-\sigma=-1/\sigma$, e ambas satisfazem a \emph{mesma}
borda --- logo $D$ e $H$ são invariantes e a involução do corpo \textbf{não move $w$}. É preciso o
segundo lado, e são \emph{dois}: o \textbf{sinal da pressão} ($w\mapsto-w$) e o \textbf{sentido do
grau} ($m\mapsto-m$). Comutam --- mas reflectem em \emph{centros diferentes}, $0$ e $-1$, porque
$w+1=2m/(3\sqrt D)$ é que é ímpar em $m$. Por isso a órbita tem \textbf{quatro}:

\begin{center}\small
\begin{tabular}{@{}lrrrr@{}}
\toprule
$m$ & $w$ & $-w$ & $w(-m)$ & $-w(-m)$ \\
\midrule
$5$ & $-1+\tfrac{10}{87}\sqrt{29}$ & $+1-\tfrac{10}{87}\sqrt{29}$ & $-1-\tfrac{10}{87}\sqrt{29}$ & $+1+\tfrac{10}{87}\sqrt{29}$ \\
$0$ (ponto fixo) & $-1$ & $+1$ & $-1$ & $+1$ \\
\bottomrule
\end{tabular}
\end{center}

\noindent E não há aproximação nenhuma nesta tabela: $w+1=\frac{2m}{3D}\sqrt D$ é um elemento
\emph{exacto} de $\mathbb{Q}[\sqrt D]$, e todas as contas se fazem no coeficiente, por produto
cruzado de inteiros. \textbf{Não se tira raiz, e por isso não há tolerância} --- como não há no
\code{MOVE}.

\smallskip
\noindent \textbf{O ponto fixo não é onde a órbita passa: é onde ela degenera} --- quatro estados
caem para dois --- e é isso que faz dele ponto fixo, aqui como no trial.

\subsection*{E $w=1/3$ não pertence à família --- e isso é resultado}

Impondo $w=\tfrac13$ em $w=\frac{2m}{3\sqrt D}-1$ com $D=m^2+4$ vem $2m=4\sqrt D$, logo
$m^2=4D=4m^2+16$, isto é $12m^2=-64$: \textbf{sem raiz real nenhuma} --- não só sem raiz inteira.
Varridos os $m$ de $-999$ a $999$: \textbf{nenhum}, e o próprio sinal negativo di-lo antes da varredura
(coerente com $w+1$ subir \emph{sempre por baixo}). Não é falha da derivação: é \emph{onde a
radiação vive} --- na fronteira, que é o que o ponto fixo em $w=-1$ já dizia. A família descreve os
corpos; a radiação pura é o limite deles.

\medido{a lei de diluição sai da continuidade \emph{com} a constante --- $\rho\,a^{3(1+w)}=C$ conservado em $36$ escalas para seis valores de $C$, sem um desvio, e a involução $C\mapsto1/C$ com \emph{dois} pontos fixos, o positivo a dar a unidade --- e os expoentes dão $a^{-4}$, $a^{-3}$ e constante para os três
conteúdos, com $w=-\tfrac13$ neutro; a órbita da bidualidade tem $4$ estados em $m=5$ e degenera em $2$ em $m=0$, com as
duas involuções a reflectirem em centros diferentes ($0$ e $-1$); $w\mapsto-w$ é o simétrico e é
involução em $17$ graus, verificado pela definição e não pela contagem; e $w+1$ sobe para
$\tfrac23$ \emph{sempre por baixo}, com a distância a valer exactamente $16/(9D)$ em $12$ graus
--- o numerador é $4\cdot4$, o discriminante a aparecer. Tudo em $\mathbb{Q}[\sqrt D]$ por produto
cruzado de inteiros: \textbf{sem raiz e sem tolerância}; e o \textbf{controlo}: com o expoente que a continuidade dá, $\rho a^{3(1+w)}$ fica
invariante em todas as escalas ($0$ falhas em $5$); trocando-o, falha em $4$ de $5$
(\code{tests/cosmologia.c}).}

\section{Os dois escuros são os dois pontos fixos}

A bidualidade tem \textbf{duas} involuções que comutam, e a órbita genérica tem quatro estados. Cada
involução tem \textbf{exactamente um} ponto fixo, e nele a órbita degenera de quatro para dois:

\begin{center}\small
\begin{tabular}{@{}llll@{}}
\toprule
a involução & fixa & o conteúdo & e é invisível \\
\midrule
$w\mapsto-w$ (a pressão) & $w=0$ & \textbf{a matéria} & à pressão \\
$m\mapsto-m$ (o grau) & $m=0$, e aí $w=-1$ & \textbf{o vácuo} & ao grau \\
\midrule
\emph{nenhuma} & --- & a radiação & \emph{é a que se vê} \\
\bottomrule
\end{tabular}
\end{center}

\noindent \textbf{Um ponto fixo não se separa do seu dual pela leitura que o fixa} --- é isso que
degenerar quer dizer. A matéria \emph{é} o seu próprio dual sob a pressão, e por isso nenhuma equação
de estado a distingue: precisa da \emph{outra} involução. Aqui ``escuro'' não é metáfora, é uma
definição: \textbf{um conteúdo que a leitura disponível não separa}.

\smallskip
\noindent E a assinatura de um conteúdo é \emph{qual} involução o fixa --- a primeira, a segunda, ou
nenhuma. \textbf{São três casos, um por conteúdo}: o trial outra vez, e não há quarto. Fixo pelas
\emph{duas} é impossível: exigiria $w=0$ e $m=0$ ao mesmo tempo, e $m=0$ dá $w=-1$.

\medido{cada involução tem \emph{um} ponto fixo e são dois distintos; a órbita degenera dos dois
lados ($4\to2$ no grau, $2\to1$ na pressão) e a troca de sinal é verificada pela definição
$w+(-w)=0$ em $17$ graus; exactamente um dos três conteúdos é invisível à pressão; os três casos da
assinatura estão ocupados, um cada, e \textbf{zero} pelas duas; e o \textbf{controlo}: em $1201$
graus nenhum satisfaz as duas, e dos três conteúdos só o vácuo é membro da família --- no próprio
ponto fixo (\code{tests/escuros.c}).}

\section{Dois universos, e a estrela entre eles}\label{sec:dois}

O buraco branco não mora cá. \textbf{A estrela não está entre dois buracos do mesmo universo ---
está entre dois \emph{universos}}, e é a interface deles. O dual do factor de escala é o seu inverso,
$a_B=1/a_A$: um contrai exactamente o que o outro expande.

\begin{center}\small
\begin{tabular}{@{}lll@{}}
\toprule
universo $A$ & expande & o buraco \textbf{negro} dissipa --- e a expansão leva \\
\textbf{a estrela} & $a=1/a$ & \textbf{a interface}: os dois sentidos, e não paga \\
universo $B$ & contrai & o buraco \textbf{branco} absorve --- e a contração traz \\
\bottomrule
\end{tabular}
\end{center}

\subsection*{Parseval, e é multiplicativo}

Com $\rho=C\,a^{-k}$ em cada um, o $a^{-k}$ de um cancela o $a^{+k}$ do outro:
\[
  \boxed{\;\rho_A\,\rho_B = C_A C_B\;}
\]
\textbf{A energia não se perde ao atravessar --- muda de domínio.} E é \emph{multiplicativo} e não
aditivo, porque o corpo é multiplicativo: é o $|N|=1$ outra vez, exactamente no sítio onde uma soma
teria sido o erro.

\subsection*{E o corpo negro sai de graça}

$\rho\sim a^{-(d+1)}$ e $T\sim a^{-1}$ --- \emph{as duas já derivadas} --- compõem-se em
\[
  \rho \sim T^{\,d+1}, \qquad\text{e em } d=3: \quad \rho\sim T^4 .
\]
Não se citou lei nenhuma: compuseram-se dois expoentes que já lá estavam. E emitir e absorver são o
par \emph{passo a passo}, não uma compensação no fim --- o que $A$ dissipa é o que $B$ recebe, em
cada passo.

\subsection*{E a antimatéria entra pela terceira involução}

São \textbf{três}, e comutam:

\begin{center}\small
\begin{tabular}{@{}lll@{}}
\toprule
a involução & fixa & separa \\
\midrule
$w\mapsto-w$ (a pressão) & $w=0$ & ocupado de vazio \\
$m\mapsto-m$ (o grau) & $m=0$ & é o vácuo \\
$a\mapsto1/a$ (\textbf{o universo}) & $a=1$ & matéria de \textbf{antimatéria} \\
\bottomrule
\end{tabular}
\end{center}

\noindent Com duas delas a órbita tem quatro, e são os quatro conteúdos:

\begin{center}\small
\begin{tabular}{@{}lll@{}}
\toprule
 & universo $A$ (expande) & universo $B$ (contrai) \\
\midrule
ocupado & matéria & \textbf{antimatéria} \\
vazio & matéria escura & \textbf{antimatéria escura} \\
\bottomrule
\end{tabular}
\end{center}

\noindent \textbf{E todas são biduais}: cada uma tem par pelo universo \emph{e} par pela pressão ---
nenhuma fica sozinha.

\medido{o expoente do corpo negro sai da composição sem dividir expoentes, dando $T^4$ em $d=3$; o
produto $\rho_A\rho_B$ vale o produto das constantes em todas as escalas; em $40$ passos o que $A$
emite é exactamente o que $B$ absorve e o total não se move --- com um lado só, cai; há \emph{um}
ponto onde $a=1/a$; os quatro conteúdos são distintos, todos com par pelas duas involuções, e elas
comutam em ordens genuinamente trocadas; e o \textbf{controlo}: com $a_B=k/a_A$ e $k\neq1$ o produto
deixa de se conservar nos quatro casos (\code{tests/parseval.c}).}

\section{A temperatura, e a entropia dos dois lados}

\subsection*{A temperatura é o inverso do raio}

A chave é uma \emph{contagem}: em p.u.\ um buraco negro tem \textbf{uma} escala independente --- o
raio, e a massa \emph{é} o raio, não uma segunda. Logo o comprimento de onda característico só pode
ser o próprio raio, e como $T$ vai com o inverso do comprimento de onda,
\[
  \boxed{\;T=\frac1r\;}
\]
\textbf{e isso \emph{é} a involução do universo} --- a mesma $a\mapsto1/a$. O buraco negro não tem lei
própria: tem a involução aplicada ao seu único comprimento. E o \textbf{ponto fixo é $r=1$}, onde a
temperatura vale a unidade: \emph{a estrela outra vez}.

\subsection*{A entropia é a área, e a evaporação sai por expoentes}

$S\sim\Omega_{d-1}r^{\,d-1}$, com o mesmo $\Omega$ que saiu por recorrência no §\ref{sec:campo2}. O
expoente é $d-1$ e não $d$: é a \textbf{codimensão}, a mesma que deu o expoente da força.
\emph{A entropia mora na fronteira, não no interior.}

\smallskip
\noindent E a evaporação: a potência é a área vezes $T^{d+1}$, e com área $\sim r^{d-1}$ e $T\sim1/r$
os expoentes somam em $-2$ --- \textbf{o mesmo em toda a dimensão}, o que não era óbvio antes de se
contar. Integrando, a vida vai com $r^3$, e como a massa \emph{é} o raio, com $M^3$.

\subsection*{E a entropia do buraco branco é a mesma lei, no raio dual}

Não há segundo princípio. Pondo o raio dual em $S=(\text{raio})^{d-1}$, o expoente vira simétrico:
\[
  S_{\text{negro}}=r^{\,d-1}\ \text{(cresce --- o \emph{máximo})},\qquad
  S_{\text{branco}}=r^{-(d-1)}\ \text{(decresce --- o \emph{mínimo})}.
\]

\noindent E daí \textbf{duas} coisas, que não são a mesma:

\begin{center}\small
\begin{tabular}{@{}ll@{}}
\toprule
$S_{\text{negro}}\cdot S_{\text{branco}}=1$ & Parseval, \textbf{multiplicativo} \\
a soma dos expoentes $=0$ & porque a entropia \emph{já é} um logaritmo: \textbf{aditivo} \\
\bottomrule
\end{tabular}
\end{center}

\noindent Contar caminhos \emph{é} somar expoentes. O par das densidades multiplica para uma
constante; \textbf{o par das entropias soma para zero}.

\subsection*{Logo a segunda lei vale de \emph{um} lado}

A entropia do negro cresce em todos os passos, a do branco decresce em todos, e \textbf{o par não se
move em nenhum}. Logo ``a entropia cresce'' é verdade e \emph{não é lei do universo --- é lei de
metade dele}. A lei do par é que a soma não muda, e \textbf{a seta do tempo não é propriedade do
tempo: é uma polaridade lida sem a outra}. No ponto fixo $r=1$ as duas coincidem, e aí a seta não
aponta para lado nenhum.

\medido{o posto das dimensões dá \emph{uma} escala independente; dobrar o raio multiplica a área por
$2^{d-1}$ e não por $2^d$; $dr/dt$ vai com $r^{-2}$ em todas as dimensões e a vida com $r^3$; a
involução $r\leftrightarrow1/r$ fecha com resíduo $0$ em $40$ raios com \emph{um} ponto fixo; o
produto das temperaturas vale a unidade em todos; o produto das entropias vale a unidade em $60$
casos e a soma dos expoentes é zero, com a ponte produto-$=$-soma verificada sem avaliar um
logaritmo; em $19$ passos o negro cresce em \emph{todos}, o branco decresce em \emph{todos} e o par
move-se em \emph{nenhum}; e o \textbf{controlo}: trocando a codimensão pelo volume, ou o inverso por
um múltiplo dele, o par deixa de fechar (\code{tests/hawking.c}, \code{tests/entropia\_dual.c}).}

\section{As constantes analíticas}

Saem de \textbf{duas} fontes, e não há terceira.

\begin{center}\small
\begin{tabular}{@{}lll@{}}
\toprule
a fonte & a constante & o que a define \\
\midrule
\textbf{o relógio} & $\pi$ & o \emph{meio-período}: gerador da Lei 2, alvo da Lei 1 \\
 & $e$ & o \emph{fluxo}: onde a soma vira produto \\
\textbf{a família} & $\varphi=\sigma_1$ & o ponto fixo da borda $\sigma^2=m\sigma+1$ \\
 & $\sigma_m$ & um por metal, e só esses \\
\bottomrule
\end{tabular}
\end{center}

\noindent \textbf{$\pi$ não pede geometria.} A Lei 2 fornece o gerador $J$, o fluxo percorre-o, e a
Lei 1 fornece o \emph{alvo}: $-1$. O meio-período é onde o fluxo lá chega, é único dentro de um
período, e a razão período/meio-período é exactamente $2$ --- \emph{a meia volta, outra vez}.

\smallskip
\noindent \textbf{$e$ não pede uma série.} $(1+1/n)^n$, em numerador e denominador inteiros, cresce
estritamente e nunca passa de $3$: monótono e limitado, o limite existe.

\smallskip
\noindent \textbf{E os $\sigma_m$ são irracionais por construção}: nenhum inteiro satisfaz
$c^2=mc+1$. É por isso que a família é uma família e não uma lista.

\smallskip
\noindent \textbf{E as três encontram-se no mesmo sítio}, que é o que faz delas um sistema: o alvo de
$\pi$ é o $-1$ da Lei 1 --- o mesmo número, não um parecido ---, o fluxo de $e$ leva a soma ao
produto, que é a involução que a constante de integração atravessa, e $\sigma_1$ é $\varphi$.

\medido{o meio-período existe, é único e a razão é $2$; $(1+1/n)^n$ cresce em $8$ passos sem falha e
nunca passa de $3$ (para em $n=9$ porque o produto cruzado deixa de caber no tipo, e um número que
não cabe mente antes de falhar); nenhum dos $328$ candidatos inteiros satisfaz a borda; $\Lambda=3/D$
em seis graus, com $3/4$ no ponto fixo; e o \textbf{controlo}: dos seis alvos testados quatro não têm
meio-período nenhum, e $\Lambda$ tem um só valor por grau (\code{tests/analiticas.c}).}

\section{A renormalização: o número é da régua, a meia volta é do objecto}\label{sec:renorm}

Renormalizar é mudar a régua sem mudar o objecto, e o relógio mostra-o num renque.

\smallskip
\noindent \textbf{A velocidade máxima é meia volta.} Num relógio de $q$ marcas a distância entre
duas é a \emph{menor} das duas voltas, $d(p)=\min(p,\,q-p)$ --- andar mais que meia volta é andar
para trás pelo outro lado. Logo
\[
  \max_p\,d(p)=\frac{q}{2},\qquad\text{atingido em } p=q-p .
\]
\textbf{O máximo é a involução}: o único ponto onde os dois sentidos dão o mesmo número, e \emph{sem
sentido definido não há órbita}. Não se postulou velocidade limite nenhuma --- ela saiu de o círculo
ser fechado.

\smallskip
\noindent \textbf{E é isso que se renormaliza.} O \emph{número} muda com a régua --- $2$, $4$, $8$,
$16$ conforme $q$ --- mas a \emph{razão} $d/q$ é $\tfrac12$ em toda a escala. E a velocidade máxima
só se \emph{atinge} em $q$ par: em $q$ ímpar não existe $p$ com $p=q-p$, e nenhuma marca lá chega.

\subsection*{Os degraus são os metais, e o fluxo preserva a forma}

A escala não sobe continuamente: sobe por degraus, e os degraus são os metais. O fluxo é
\[
  x\;\longmapsto\;m+\frac1x,\qquad\text{com ponto fixo}\quad \sigma_m=\frac{m+\sqrt{m^2+4}}{2}.
\]
E o que faz dele uma renormalização é o que ele \textbf{preserva}: ao longo de todo o fluxo a forma
da borda
\[
  Q(n,d)=n^2-m\,n\,d-d^2
\]
só muda de \emph{sinal} --- o módulo fica. O número $n/d$ muda a cada passo; a classe não. E a
semente não escolhe um valor arbitrário: escolhe \emph{qual órbita} da forma se percorre.

\smallskip
\noindent \textbf{E a taxa é a própria régua.} O denominador obedece à recorrência da borda,
$d_{k+1}=m\,d_k+d_{k-1}$, logo cresce por $\sigma_m$ a cada passo --- e o erro cai pelo quadrado
disso. \emph{O ritmo com que a espiral se aproxima da sua régua é a régua.}

\smallskip
\noindent E há uma dimensão onde as duas réguas comensuram: $\sigma_4=\varphi^3$, exactamente. Medir
com a régua do ouro e com a régua própria dá o mesmo número a menos de um expoente, e isso não
acontece em quase nenhuma outra.

\medido{a razão $d/q$ é $\tfrac12$ em sete escalas sem desvio; em $30$ escalas pares há
\emph{exactamente uma} involução em cada e em $29$ ímpares não há nenhuma; o fluxo preserva $|Q|$ em
cinco metais e duas sementes sem se mover um único passo (de $m/1$ sai $1$, de $1/1$ sai $m$); o
denominador cresce pela recorrência da borda em $120$ passos com $0$ falhas, com a razão sempre
entre $m$ e $m+1$, que é onde a borda põe a raiz; $\varphi^3=1+2\varphi=\sigma_4$ calculado em
$\mathbb{Z}[\varphi]$ \emph{sem avaliar uma raiz}; e o \textbf{controlo}: nenhum inteiro fora do
próprio metal satisfaz $c^2=mc+1$, falhando os $25$ candidatos (\code{tests/renormaliza.c}).}

\section{Vestir a roupa}\label{sec:roupa}

Sair do por-unidade não é ganhar acesso ao absoluto: é passar a ter uma base que o objecto forneceu.
E a ordem importa --- \textbf{a estrutura deriva-se primeiro, e as constantes vêm no fim}. Se a
derivação precisasse do valor para chegar ao valor, não era derivação.

\subsection*{A velocidade máxima, e o $32$}

Se a velocidade máxima é meia volta, a distância máxima \emph{a essa velocidade} é o percurso que usa
\textbf{todas} as arestas e regressa. No hipercubo de dimensão $n$ há $n\,2^{n-1}$ arestas ---
\textbf{$32$ em $n=4$} --- e esse percurso existe exactamente quando todo o grau é par. \emph{O grau
do hipercubo é a dimensão}, logo fecha nas pares e não fecha nas ímpares. Nada disto pediu constante
nenhuma.

\subsection*{Quantas escalas tem a roupa}

\textbf{Três}, e o número não foi escolhido: é o posto das dimensões das grandezas do sistema,
calculado por eliminação inteira. É também quantas entradas tem a assinatura de um corpo, e pela
mesma razão --- \emph{o trial não tem quarto estado}. Fixadas as três, toda a grandeza fica com o seu
expoente e não sobra nada por dizer.

\subsection*{O que sai, e o que não sai}

\begin{center}\small
\begin{tabular}{@{}lll@{}}
\toprule
 & sob troca de roupa & donde \\
\midrule
constante \textbf{dimensional} & muda pelo factor dos seus expoentes & \textbf{sai da régua} \\
número \textbf{adimensional} & \emph{não se move} & não sai da régua \\
incerteza \textbf{relativa} & é razão: \emph{não se move} & nenhuma régua a remove \\
\bottomrule
\end{tabular}
\end{center}

\noindent Fixadas as escalas, uma grandeza de expoentes $(a,b,c)$ vale $L^aT^bM^c$ e não tem
liberdade nenhuma. Mas um número puro vale o mesmo em qualquer roupa --- \textbf{e é aí que a
derivação para}. É importante que se veja onde.

\subsection*{E a ordem: o $8\pi$ primeiro, depois o passo, depois a régua}

Ao sair do por-unidade, \textbf{o primeiro que aparece não é uma escala}. Um número puro não depende
da roupa --- está lá \emph{antes} de haver escala nenhuma --- e por isso vem antes. E o primeiro é o
$8\pi$, que saiu de contar: a esfera e o traço.

\begin{center}\small
\begin{tabular}{@{}rlll@{}}
\toprule
 & o que entra & de onde & o que fornece \\
\midrule
$1$ & \textbf{o número puro} & contar: a esfera e o traço & $8\pi$, e não se move em roupa nenhuma \\
$2$ & \textbf{o passo} & o \emph{relógio} & período $4$; meia volta é $\pi$ \\
$3$ & \textbf{a régua} & a \emph{família metálica} & $\sigma$ é a borda, e $[v]=[\sigma]$ \\
\bottomrule
\end{tabular}
\end{center}

\noindent \textbf{E com os três sai toda a constante dimensional}: cada uma é
$(\text{passo})^a(\text{régua})^b$, e o expoente é o seu --- trocada a roupa, move-se pelo factor que
os \emph{seus} expoentes obrigam e por mais nenhum. \emph{São três entradas e não há quarta}, e
nenhuma delas vem de fora: o puro vem de contar, o passo vem do relógio, a régua vem da borda.

\smallskip
\noindent E cada uma é indispensável: sem o passo, quatro das seis grandezas ficam indeterminadas;
sem a régua, cinco; e sem o número puro a equação de campo não tem coeficiente.

\smallskip
\noindent \textbf{O que fica de fora fica à vista.} Um número puro que \emph{não} venha de uma
contagem --- a constante de estrutura fina é o caso --- não sai daqui, e não é por falta de escala:
é por não haver o que contar. \emph{O que está fora do alcance é o que não se conta}, e essa é a
fronteira, dita sem rodeio.

\subsection*{E as constantes, com nome}

O por-unidade não é notação: é um \textbf{corte}. Fora dele há três dimensões independentes; dentro,
as duas relações --- velocidade $=1$ e acoplamento $=1$ --- cortam duas, e \textbf{sobra uma}. E a
régua $\sigma$ não é a segunda escala: é uma \emph{densidade}, logo uma razão, logo um número puro.
Daí a tabela fica curta:
\[
  \boxed{\;\text{cada constante} \;=\; (\text{passo})^{a} \times (\text{número puro})\;}
\]

\begin{center}\small
\setlength\tabcolsep{5pt}
\begin{tabular}{@{}lll@{}}
\toprule
 & o coeficiente & o que é \\
\midrule
$c$ & $\sigma$ & \textbf{é a régua} --- o máximo \emph{é} $\sigma$, não um limite de fora \\
$G$ & $8\pi$ & o acoplamento: a esfera \emph{e} o traço \\
$k$ (Coulomb) & $4\pi$ & \textbf{a área da esfera} --- o fluxo total \\
$k_B$ & $1$ & \textbf{um câmbio}, não uma constante da natureza \\
\bottomrule
\end{tabular}
\end{center}

\noindent \textbf{E daí um resultado que cai de graça}: Coulomb leva $4\pi$ e Einstein leva $8\pi$, e
a razão entre eles é \emph{exactamente} $2$ --- o factor do traço, derivado em §\ref{sec:campo2}.

\smallskip
\noindent E $k_B$ não é constante da natureza: é um \textbf{câmbio}. Dizer que a temperatura é
energia é escolher medi-las na mesma unidade --- pô-lo em $1$ não é aproximação, é \emph{reconhecer
que são a mesma grandeza}.

\subsection*{E as restantes, todas em função de $c$}

Há \emph{uma} amarra que prende o lado electromagnético inteiro:
\[
  \boxed{\;\mu\varepsilon=\frac1{c^2}\;}
\]
e não é coincidência de unidades: é a afirmação de que a perturbação se propaga à \textbf{velocidade
máxima}. Impor $v=c$ em $v^2=1/(\mu\varepsilon)$ deixa \emph{uma} solução --- e em p.u.\ é trivial,
porque $c=1$.

\begin{center}\small
\setlength\tabcolsep{5pt}
\begin{tabular}{@{}lll@{}}
\toprule
 & em $c$ & o que é \\
\midrule
$\mu\varepsilon$ & $c^{-2}$ & \textbf{a amarra} --- a perturbação vai à velocidade máxima \\
$Z$ (impedância) & $c^{+1}$ & $\mu c$ --- o vácuo tem impedância \\
$k$ (Coulomb) & $c^{+2}\cdot4\pi$ & $\mu c^2/4\pi$ --- a esfera outra vez \\
$E$ de repouso & $c^{+2}$ & $mc^2$: a massa \emph{é} energia, pela régua ao quadrado \\
$a$ (radiação) & $c^{-1}$ & $4\sigma_{\!S\!B}/c$ --- o $c$ entra a dividir \\
$\sigma_{\!S\!B}$ & $c^{0}$ & o câmbio $k_B$ fixa-a; não traz escala nova \\
\bottomrule
\end{tabular}
\end{center}

\noindent \textbf{Cada uma é uma potência de $c$ vezes um puro}, e o puro é sempre um dos que já se
contaram. Do lado térmico a amarra é o \emph{câmbio}: fixado $k_B$, a radiação segue.

\smallskip
\noindent E a \textbf{carga não entra}: $e^2$ pede $\alpha$, que é um puro \emph{sem contagem por
trás}. Não fica de fora por falta de escala --- fica por não haver o que contar.

\subsection*{E o $c^2$ sai da espiral --- não é uma potência acrescentada}

Em p.u.\ $c=1$, logo \textbf{$E=m$}: a energia \emph{é} a massa, não uma quantidade proporcional a
ela. Então o $c^2$ de $E=mc^2$ tem de vir de algum lado --- e vem da \textbf{borda}. Como $c=\sigma$
e $\sigma^2=k\sigma+1$,
\[
  \boxed{\;\sigma^{\,j} = F_j\,\sigma + F_{j-1}\;}
\]
com $F$ a recorrência do metal. \textbf{O quadrado da régua é linear na régua}, e o mesmo vale para
\emph{toda} a potência: a espiral inteira cabe em \textbf{duas} coordenadas.

\smallskip
\noindent E é por isso que a torre fecha em dois: a borda é de grau dois, e não há terceiro
coeficiente a guardar. $c^2$ não acrescenta dimensão nenhuma --- é $c$ mais uma unidade, a menos do
grau do metal.

\subsection*{E o que é a massa: o segundo momento, e o grau decide}

``Quantidade de matéria'' é circular. O que o sistema dá vem da \textbf{cruz}, que já tem duas
coordenadas: com $x^{\dagger}=2c-x$, o produto é $x(2c-x)=2cx-x^2$, donde
\[
  n\,c^2 - \textstyle\sum x\,x^{\dagger} \;=\; \sum (x-c)^2 ,
\]
exactamente. \textbf{A primeira coordenada da cruz é o primeiro momento --- o centro, o que não se
moveu --- e a segunda é o segundo.} Logo a \textbf{massa é momento quadrático}, e não por analogia:
pela mesma conta.

\subsection*{E a massa é vectorial porque é \emph{ordenada}}

\textbf{No bidual sempre existe ordem} --- e uma ordem sobre $n$ partes precisa de $n$ componentes
para ser carregada. Um escalar não a carrega: \emph{colapsa-a}. As $120$ ordens de cinco partes dão
\textbf{uma} só soma, e o número perde $119$.

\smallskip
\noindent \textbf{E na física é o mesmo}: cada corpo tem uma ordem em cada partícula, \emph{senão
ninguém crescia nem se curava}. Crescer e curar são reconstrução \textbf{ordenada} --- cada parte tem
de saber onde vai --- e se a massa fosse um número essa informação não estava em lado nenhum.

\smallskip
\noindent Ser corpo é o que lhe dá as $n$-uplas onde a ordem cabe; a bidualidade é o que garante que
a ordem existe. Em $\mathbb{R}^n$ a segunda coordenada da cruz não é um número: é a \emph{matriz}
\[
  M_{ij} = \sum (x_i-c_i)(x_j-c_j) ,
\]
e a massa escalar é apenas o seu \textbf{traço}. Logo \emph{a massa tem direcção} --- um corpo pode
ter massa diferente em eixos diferentes, e o número não o vê --- e a bidualidade é o que impede o
colapso, embora seja a consequência e não a razão. Mede-se pelo que se perde: de $168$
distribuições o tensor separa $84$ e o traço apenas $26$ --- \textbf{$58$ distinções desaparecem} ao
ficar só com o escalar. É a contagem de sempre: dividir perde, e o que se perde é o segundo lado.

\smallskip
\noindent \textbf{E o grau decide qual é qual:}

\begin{center}\small
\begin{tabular}{@{}llll@{}}
\toprule
grau & raízes & há par? & o que dá \\
\midrule
$1$ & uma & não & transporte, \textbf{momento linear}, meia volta \\
$2$ & duas & \textbf{sim} & o par, a \textbf{massa}, a volta fechada \\
\bottomrule
\end{tabular}
\end{center}

\noindent Uma equação do primeiro grau resolve-se e acabou --- é por isso que converter uma sequência
noutra sai barato. Uma do segundo tem \emph{duas} raízes, com produto $-1$ --- que é a Lei 1 --- e é o
par que faz a volta fechar. \textbf{E a espiral cabe em duas coordenadas porque a borda é de grau
dois}: não são dois factos, é um.

\smallskip
\noindent Daí cai, sem se pedir, que \emph{o que anda à velocidade máxima não tem massa}: à velocidade
máxima ir e voltar são o \emph{mesmo} caminho --- o único ponto onde $p=q-p$ --- logo não há volta
distinta a fechar, nem dispersão própria. Não é que a massa seja zero por acidente: é que a distinção
entre ida e volta desapareceu.

\subsection*{E porque é o quadrado: a dimensão não decide, a volta decide}

Como $\sigma$ é um número \emph{puro}, $mc$ e $mc^2$ têm a \textbf{mesma dimensão} neste sistema --- o
argumento de que $c^2$ traria velocidade ao quadrado não se aplica: aqui a velocidade não é uma
escala, é uma densidade. O que os separa é
a \textbf{contagem de passagens} pela régua, e a resposta vem de onde vem tudo o resto: da reversão.

\begin{center}\small
\begin{tabular}{@{}lll@{}}
\toprule
$mc$ & \emph{uma} passagem --- não reverte & é o \textbf{momento} \\
$mc^2$ & ida \emph{e} volta --- resíduo $0$ & é a \textbf{energia} \\
\bottomrule
\end{tabular}
\end{center}

\noindent \textbf{Logo $E=mc^2$, e pela volta.} O que sobrevive é o que fecha, e meia volta não
fecha.

\medido{a dimensão \emph{não} separa $mc$ de $mc^2$ nas quatro roupas, porque $\sigma$ é puro; uma
passagem pela régua deixa resíduo em $220$ casos e a ida-e-volta deixa \emph{zero}; das $32$
velocidades do relógio, $31$ fecham órbita e a \emph{única} que não fecha é exactamente a máxima; a
identidade $nc^2-\sum xx^\dagger=\sum(x-c)^2$ fecha em $99$ casos sem desvio; e o tensor separa $84$
distribuições onde o traço separa $26$ --- $58$ distinções perdidas no escalar; $E=m$ em p.u.\ com resíduo $0$ em $41$ massas; $\sigma^j=F_j\sigma+F_{j-1}$ em $72$ potências
de seis metais sem um desvio, com o coeficiente de $\sigma$ a ser \emph{exactamente} o Fibonacci do
metal --- verificado em $\mathbb{Z}[\sigma]$, sem avaliar uma raiz; e quatro mutações na recorrência,
quatro acusam.}

\medido{impor $v=c$ deixa \emph{uma} solução para $\mu\varepsilon$ em $11$ velocidades; e cada
expoente da tabela é \emph{derivado} e não escrito --- o de $\mu\varepsilon$ é menos duas vezes o da
velocidade, o de $E$ é duas vezes, o de $a$ é o simétrico, e o coeficiente de Coulomb tem de bater
com o $4\pi$ medido três secções acima. Seis mutações na tabela, seis acusam
(\code{tests/constantes.c}).}

\medido{o posto é $3$ fora do p.u.\ e $1$ dentro, com as duas relações a cortarem exactamente duas
dimensões; $\sigma$ não se move em roupa nenhuma e o passo move-se em todas menos a trivial; as oito
entradas da tabela têm todas origem --- quatro levam o passo, quatro são só puros, \emph{nenhuma} sem
origem; a razão Einstein/Coulomb bate com o factor do traço; e o \textbf{controlo}: com posto $3$ a
tabela precisaria de $24$ números e com posto $1$ precisa de $8$, e todos os puros vêm de uma
contagem --- $\pi$ do relógio, $8\pi$ da esfera e do traço, $\sigma$ da borda
(\code{tests/constantes.c}).}

\medido{as arestas dão $32$ em $n=4$ e o percurso fecha nas dimensões pares e não nas ímpares,
nomeadas uma a uma; o posto das dimensões de nove grandezas é $3$, com os expoentes de $G$
\emph{derivados} de $G=Fr^2/m^2$ e não escritos; a razão adimensional dá o mesmo número nas cinco
roupas enquanto a velocidade e $G$ mudam em todas; o $8\pi$ fecha por recorrência sem avaliar $\pi$ e
não se move em quatro roupas; o relógio dá período $4$ e meia volta $\pi$; nenhum dos $328$
candidatos inteiros satisfaz a borda; as seis grandezas movem-se pelo factor dos seus próprios
expoentes, $0$ desvios; e o \textbf{controlo}: tirando o passo quatro ficam indeterminadas, tirando
a régua cinco --- nenhuma das três entradas é dispensável, e nenhuma vem de fora
(\code{tests/vestir.c}, \code{tests/sair\_pu.c}).}

\section{O corpo estelar completo: a torre nos lados, a plena no meio}\label{sec:completo}

Até aqui a trindade foi lida nos sinais e nas densidades --- em $\mathbb{R}$, em grau baixo. Falta
dizer \emph{de que grau} é cada um dos três, e a resposta fecha o sistema, porque não é uma escolha:
\textbf{deriva-se}. E derivá-la é o salto que a teoria devia --- \emph{toda ela se construiu até aos
complexos}, o plano, o grau dois; os dois buracos e a estrela pedem um andar acima.

\subsection*{A interface é reversível, logo é a plena --- grau seis}

\begin{teorema}[a estrela tem grau seis]\label{thm:estrelaseis}
Uma interface reversível --- em que \emph{ler é escrever ao contrário} --- só existe no grau em que a
soma e o produto coincidem, e esse grau é \textbf{um só}: o seis.
\end{teorema}

\begin{proof}
Reversível quer dizer que a leitura é a escrita com o sinal trocado: uma operação, dois sentidos
($\texttt{MOVE}(\cdot,\pm1)$). Para que ler e escrever sejam \emph{a mesma} operação --- e não duas
que por acaso se desfazem --- é preciso que pôr e ler não se distingam, isto é, que a soma e o
produto deem o mesmo sobre as mesmas partes. Isso vale num grau só, o \textbf{seis}, o único inteiro
com $1+2+3=6=1\cdot2\cdot3$ (a \emph{plena}): ali as duas operações coincidem e \emph{nada distingue
o que se põe do que se lê}. Fora dele, somar não é multiplicar, escrever não é ler, e uma interface
que os confundisse dissiparia. Logo a estrela --- a única reversível --- tem grau seis, por
necessidade e não por escolha.
\end{proof}

\noindent \textbf{E é por isso que a estrela paga zero.} Não porque seja eficiente: porque no grau
seis a escrita \emph{é} a leitura, e uma operação que é a sua própria inversa não apaga bit nenhum
(§\ref{sec:medir}). A plena é o único andar onde a involução e a evolução são a mesma coisa --- e a
estrela mora exactamente aí.

\noindent \textbf{Grau seis é a hexal; o seu bit é a pental.} Na progressão das seis leis (\emph{Catálogo},
\code{obs:seis-leis}) o grau seis é a \textbf{HEXAL} (dim $6$): a \emph{interface} reversível, o relógio.
E o seu ponto fixo $x^2=-1$ --- o \textbf{bit} $i$ --- é a \textbf{PENTAL} (dim $5$, o corolário). As seis
leis não substituem a Lei 1 e a Lei 2 deste paper: \emph{estendem-nas} --- a Lei 1 é o \textbf{dual}
(dim $2$, a involução), a Lei 2 é o gerador da descida cujo ponto fixo $i$ é a pental ---, e são as
\emph{primitivas operacionais} (cada lei um operador certificado a fechar por um medidor, resíduo $0$)
que fundamentam a construção.

\subsection*{Os dois lados são a torre --- e a torre não tem topo}

Os dois buracos são duais, $a\mapsto1/a$, e cada um \emph{opera} --- soma, multiplica, inverte. Cada
um é \textbf{grau quatro} --- os quaterniões $\mathbb{H}$, o tetral, onde o produto deixa de comutar
(é preciso a ordem: qual absorve primeiro) mas ainda associa. E o salto seguinte dá \textbf{grau
oito}: $\mathbb{O}=\mathbb{H}\times\mathbb{H}$, \emph{dois grau quatro} ligados pela mesma involução.
\emph{Um lado é dois quaterniões duais.}

\smallskip
\noindent \textbf{E aqui é preciso não trazer um morto por régua.} A leitura clássica diz que a
torre \emph{pára} em oito: em dezasseis a norma multiplicativa cai. \emph{É verdade --- para um
ramo, e um só.} O teorema que o diz (Hurwitz) \textbf{classifica} as álgebras de composição
\emph{bilineares}, e não proíbe coisa nenhuma. E a saída não é forçar a régua da esfera --- a norma
$\|x\|=\sqrt{\sum x_i^2}$ --- para dentro dos inteiros: isso é medir a esfera com o compasso do
círculo, e o $\sqrt{\ }$ que aparece é o \emph{preço da régua que se trouxe}. A saída é a
\textbf{dualidade}, e ela é o teorema central deste sistema.

\subsection*{Primeiro, o que é dual}\label{sub:defdual}

\begin{definicao}[bijeção dual --- a estrela]\label{def:bijdual}
Uma \textbf{bijeção dual} entre dois objectos $A$ e $B$ é a \emph{interface reversível}: uma
correspondência $\nu$ que é o seu próprio inverso, $\nu\circ\nu=\mathrm{id}$ com resíduo $0$; que
\textbf{liga sem fundir} --- preserva os dois lados em vez de os colar num terceiro; e que \textbf{não
apaga um bit}. É a \textbf{estrela}: o ponto fixo da involução que troca os dois, a plena
(Teorema~\ref{thm:estrelaseis}, grau seis, onde soma $=$ produto), o único regime com os dois sentidos.
\end{definicao}

\begin{definicao}[dual]\label{def:dual}
$A$ e $B$ são \textbf{duais} quando existe uma bijeção dual entre eles. Nenhum é prévio: a bijeção é o
seu próprio inverso, e o que $A$ é para $B$, $B$ é para $A$. \emph{Não há um melhor} --- dá no mesmo
pela cruz (ambos leem a mesma coisa) e desce pela estaca (a troca de sinal).
\end{definicao}

\subsection*{O teorema central: Gentil e Hurwitz são duais}\label{sub:central}

\begin{teorema}[Gentil e Hurwitz são duais, pela medida]\label{thm:central}
Sejam o lado \textbf{discreto} --- \emph{Hurwitz}: a norma euclidiana $N(x)=\sum x_i^2$ (a esfera),
lida pela cruz, multiplicativa para o produto \emph{bilinear} exactamente nos graus $1,2,4,8$ --- e o
lado \textbf{contínuo} --- \emph{Gentil}: a norma $\|x\|$ e a fusão homogénea (a hipérbole), sem limite
de grau. Então os dois são \textbf{duais}, e a bijeção dual é a \textbf{estrela}, realizada pela
\emph{teoria da medida}: o integral de Lebesgue --- a soma reversível
\[
  \int_a^b f \;+\; \int_{f(a)}^{f(b)} f^{-1} \;=\; b\,f(b)-a\,f(a)
\]
--- é a correspondência que leva \textbf{contar} (o discreto, $\sum$, Hurwitz) a \textbf{integrar} (o
contínuo, $\int$, Gentil) e a traz de volta. A bijeção é $f:t\mapsto t^2$ (contar $\leftrightarrow$
raiz), e \emph{nunca se avalia uma raiz}: os dois cortes --- o do domínio (Riemann, Hurwitz) e o da
imagem (Lebesgue, Gentil) --- dão a mesma contagem.
\end{teorema}

\begin{proof}
A bijeção dual é a involução $f\leftrightarrow f^{-1}$, que é a Lei~1. A conservação acima
--- a mesma já medida nos metais --- reparte o rectângulo pelos dois cortes \emph{sem resto}, logo a
correspondência é reversível ($\nu\circ\nu=\mathrm{id}$, resíduo $0$) e é uma bijeção; é a estrela
(Definição~\ref{def:bijdual}). \emph{Contar} (a soma no domínio, o discreto, Hurwitz) e \emph{integrar}
(o integral na imagem, o contínuo, Gentil) são os dois lados desta bijeção, logo $A$ e $B$ são duais
(Definição~\ref{def:dual}). E o \textbf{limite no grau oito é do lado discreto/bilinear} --- Hurwitz
classifica o bilinear ---, \emph{não do objecto}: o lado contínuo, alcançado pela soma reversível e não
por uma raiz, não o tem.
\end{proof}

\begin{definicao}[o octonião dual --- o grau oito reversível]\label{def:octoniao-dual}
O grau oito é $\mathbb{O}=\mathbb{H}\times\mathbb{H}^{*}$: \emph{dois tecidos grau quatro} (quaterniões)
ligados pela \textbf{estrela} (a bijeção dual, Definição~\ref{def:bijdual}). Ele lê-se de \emph{dois}
modos --- e são o par do Teorema~\ref{thm:central}:
\begin{itemize}
\item pela \textbf{norma bilinear} (Hurwitz, o lado discreto) é o \emph{octonião clássico}, e aí
\textbf{perde a associatividade} --- é o preço do grau oito \emph{do lado da norma};
\item pela \textbf{dualidade} (a cruz, a estrela, a bijeção reversível --- Gentil, o lado contínuo) é o
\textbf{octonião dual}: $\nu\circ\nu=\mathrm{id}$ com resíduo $0$, a estrela \emph{liga sem fundir}, e
\textbf{não perde nada} --- a volta fecha.
\end{itemize}
É o \emph{mesmo espaço}; muda a régua. E a régua dual é reversível \emph{onde} a bilinear cai --- porque
a estrela não é bilinear, e Hurwitz só classifica as bilineares: o limite do grau oito é do lado da
norma, \textbf{não do objecto}. É o octonião que este sistema usa em toda a parte --- a interface que
liga dois corpos sem os fundir, e não dissipa.
\end{definicao}

\medido{o octonião dual realiza-se, medido, no circuito do tradutor: dois tetrais (o corpo, grau $4$)
ligados pela interface (a hexal), em dimensão $8$ \textbf{reversível} --- a volta
$\text{tex}\to\text{pdf}\to\text{tex}$ fecha com resíduo $0$ por ser inteira, e um passo com perda (a
norma que arredonda) quebra-a ($1968$ de $3000$). A \emph{reversibilidade} é o que o separa do octonião
clássico, que no grau oito perde a associatividade (\code{tests/circuito\_tradutor.c}).}

\setcounter{lei}{6}
\begin{lei}[o octonião dual --- ligar sem fundir]\label{lei:octoniao}
$\mathbb{O}=\mathbb{H}\times\mathbb{H}^{*}$, dois tecidos ligados pela estrela: $\nu\circ\nu=\mathrm{id}$,
resíduo $0$. Pela norma perde a associatividade; pela dualidade não perde nada.
\end{lei}

\noindent \textbf{É a sétima lei, e é operacional como as outras.} As seis descem $6\to1$ \emph{dentro} de
um corpo (Teorema~\ref{thm:estrelaseis}); a sétima \emph{liga dois}. O seu \textbf{operador} é a estrela
--- a bijeção dual (Def.~\ref{def:bijdual}), uma operação em dois sentidos ($\texttt{MOVE}(\cdot,\pm1)$)
---, e o seu \textbf{medidor} certifica-a a fechar por resíduo $0$: o circuito do tradutor --- dois
tetrais ligados pela interface, a volta $\text{tex}\to\text{pdf}\to\text{tex}$ exacta por ser inteira, e um
passo com perda (a norma que arredonda) quebra-a, $1968$ de $3000$ (\code{tests/circuito\_tradutor.c}). E
\emph{a interface é a hexal}: a costura não é o oito, é o seis --- o oito é o que se liga, o seis é onde
se liga. É o octonião que este sistema usa em toda a parte, a interface que não dissipa.

\subsection*{A torre dos tecidos: cada andar é o dual do de baixo}\label{sub:tecidos}

\textbf{A estrela não liga só dois corpos dados: ela \emph{gera} a torre inteira, um andar de cada
vez.} O primeiro tecido é a \textbf{contagem natural} $\mathbb{N}$ --- o discreto puro, o lado de Hurwitz no
seu grau mínimo. Todo andar acima sai do de baixo pela \emph{mesma} operação: o dual soma-se ao
original.

\begin{definicao}[tecido, e o passo de indução]\label{def:tecido}
Um \textbf{tecido} é um andar da torre dos sistemas numéricos. O passo que leva um andar ao seguinte é
\[
  T_{k+1} \;=\; T_k \;+\; T_k^{*},
\]
o par $(A,A^*)$ tomado a menos da \textbf{estaca} --- a involução $\nu$ com $\nu\circ\nu=\mathrm{id}$ e
ponto fixo \emph{na fronteira} --- e lido pela \textbf{cruz}, a medida que se conserva no passo. É a
bijeção dual (Definição~\ref{def:bijdual}) usada como \emph{construtor}.
\end{definicao}

\begin{teorema}[a torre dos tecidos é a estrela iterada; o salto $\mathbb{Q}\to\mathbb{R}$ é o de Hurwitz para Gentil]\label{thm:tecidos}
Partindo de $\mathbb{N}$ e aplicando o passo, a torre é
\[
  \mathbb{N} \;\xrightarrow{\;+\mathbb{N}^{*}\;}\; \mathbb{Z} \;\xrightarrow{\;+\mathbb{Z}^{*}\;}\; \mathbb{Q}
     \;\xrightarrow{\;+\mathbb{Q}^{*}\;}\; \mathbb{R} \;\xrightarrow{\;+\mathbb{R}^{*}\;}\; \mathbb{C} \;\longrightarrow\;\cdots,
\]
e três coisas, que são uma:
\begin{enumerate}
\item \textbf{é uma só operação, não quatro.} Classicamente são quatro construções distintas --- as
  diferenças ($\mathbb{Z}$), as razões ($\mathbb{Q}$), a completação ($\mathbb{R}$), os pares algébricos ($\mathbb{C}$). Aqui são o
  \emph{mesmo} passo $T{+}T^{*}$: a estaca dá a involução, a cruz conserva a medida, e o ponto fixo na
  fronteira é o que fica de \emph{fora} do par (o $0$ da diferença, o corte da completação).
\item \textbf{o salto qualitativo está num degrau só, e é o Teorema~\ref{thm:central}.} De $\mathbb{N}$ a $\mathbb{Q}$
  a estaca é \emph{algébrica} (opera, mas não completa); o degrau $\mathbb{Q}\to\mathbb{R}$ é onde \textbf{contar} vira
  \textbf{integrar} --- o discreto (Hurwitz) passa a contínuo (Gentil) pela soma reversível de Lebesgue.
  A torre não muda de mecanismo; muda de \emph{régua} nesse degrau, e é o único.
\item \textbf{a interface entre andares é uma dobra.} O bit atravessa de um tecido ao seguinte pelo
  \textbf{ponto fixo} da estrela, na dobra $\Delta=n^2+4$ do corpo (Teorema~\ref{thm:atravessa}) --- a
  estaca sobe, a cruz conserva, e a estrela fecha com resíduo $0$.
\end{enumerate}
\emph{E a indução não pára por dentro:} não há andar último, porque a régua é infinita
(Teorema~\ref{thm:atravessa}). O que a faz parar é \emph{externo} --- a \textbf{finitude do corpo} que
chega: enquanto o corpo não pára de chegar, a estaca não pára de subir e a estrela não pára de reverter.
\end{teorema}

\begin{proof}
Cada passo é a bijeção dual (Definição~\ref{def:bijdual}) instanciada: a involução $\nu$ é a diferença
($\mathbb{Z}$), a razão ($\mathbb{Q}$), o corte/limite ($\mathbb{R}$), a conjugação ($\mathbb{C}$); em todas $\nu\circ\nu=\mathrm{id}$
com resíduo $0$ e ponto fixo na fronteira, e a medida (a cruz) reparte sem resto --- é a conservação do
Teorema~\ref{thm:central}, a mesma já medida. Logo o passo está bem definido e é reversível em todos os
andares. O degrau $\mathbb{Q}\to\mathbb{R}$ é o único em que a estaca deixa de ser algébrica: $\mathbb{Q}$ opera mas não é
completo, e a completação é a soma reversível $\int f+\int f^{-1}=b\,f(b)-a\,f(a)$ (Teorema~\ref{thm:central}),
isto é, o cruzamento de Hurwitz para Gentil. Que a torre não tenha topo é o crescimento sem tecto da
dobra $n^2+4$ (Teorema~\ref{thm:atravessa}); a paragem, quando há, é a finitude do corpo, não da
régua. A construção completa, com $\mathbb{N},\mathbb{Z},\mathbb{Q},\mathbb{R}$ levantados um do outro pela estaca e conservados pela
cruz, é a do \emph{Catálogo}.
\end{proof}

\noindent \textbf{E é por isto que $\mathbb{N}$ é o primeiro tecido, não um pré-requisito.} A contagem natural
não está \emph{antes} da teoria: é o andar zero dela --- o lado de Hurwitz no grau mínimo ---, e tudo o
mais é a estrela aplicada a si mesma. \emph{Um tecido, o dual, e a subida: a torre inteira sai daí.}

\subsection*{O que gera os dois: um único bit}\label{sub:bit}

Falta dizer \emph{de que} os dois lados são feitos. E a resposta é operacional, não física: \textbf{um
único bit}. A \emph{Estrela} deste sistema é a \textbf{Estrela algébrica} --- a interface de involução
com resíduo zero ---, e não uma estrela do céu; daqui por diante escreve-se só \emph{Estrela}, e nada
nela é cosmologia.

\begin{definicao}[bit operacional e suporte neutro]\label{def:bit}
A estrutura dos slots é regida pela presença ou ausência de uma única unidade binária $b\in\{0,1\}$ sob
operadores reversíveis de \emph{convolução} e \emph{deconvolução}:
\begin{itemize}
\item $b=1$: \textbf{presença} --- o slot activo no reticulado discreto (a estrutura visível);
\item $b=0$: \textbf{ausência} --- o \textbf{suporte neutro}, o elemento neutro (o ``vazio'', apenas
como elemento neutro, sem física).
\end{itemize}
\end{definicao}

\begin{teorema}[um único bit desenha o discreto e o contínuo, e o dual é a sua ausência]\label{thm:bit}
Basta um único bit $b\in\{0,1\}$ sob o par reversível \textbf{convolução} (o gato, $\times\sigma$,
\emph{sobe}) e \textbf{deconvolução} (o esquilo, $\times\sigma'=\div\sigma$, \emph{desce}) para gerar os
dois lados do Teorema~\ref{thm:central}: o \textbf{discreto} é o próprio bit (a marca no reticulado); o
\textbf{contínuo} é o rasto da convolução (o fluxo, cuja razão $\to\sigma$); e o \textbf{dual} é a
\emph{ausência} do bit ($b=0$, o suporte neutro). E fecha: $\text{gato}\circ\text{esquilo}=\mathrm{id}$,
resíduo $0$ --- a Estrela, $(A\!\times\!b)\circ(A^{-1}\!\times\!b')=\mathrm{id}$.
\end{teorema}

\begin{proof}
Semeie-se um único bit --- $b=1$ num slot, $b=0$ (neutro) nos restantes: o \emph{delta} $[1,0]$. A
convolução ($\times\sigma$, a matriz companheira $A_m$) aplicada ao delta \textbf{desenha} o reticulado
métrico --- em $m=1$ dá $[1,0],[1,1],[2,1],[3,2],[5,3],\dots$, os Fibonacci ---, e a razão dos seus
termos converge a $\sigma$: o contínuo é o rasto do discreto. A deconvolução ($\times\sigma'$, $A^{-1}$)
devolve o bit único, resíduo $0$: a interface reverte. O dual não se constrói: é a ausência, $b=0$, dada
de graça. Logo \emph{um} bit gera o discreto (ele próprio) e o contínuo (a sua convolução), e o dual é o
bit que falta. \emph{É por isto que discretização e continuidade coexistem} --- não são dois objectos,
são um bit e o seu rasto.
\end{proof}

\medido{o bit único $[1,0]$ sob o gato desenha o metal --- Fibonacci em $m=1$, a razão $\to\varphi$ ao
oitavo passo ---, e o esquilo devolve-o com resíduo $0$; a volta $\text{gato}\circ\text{esquilo}=\mathrm{id}$
fecha em $\mathbb{R}^2$ a $\mathbb{R}^8$, e o dual é a ausência $b=0$ (\code{tests/neuronio.c}).}

\subsection*{E a bijeção dual é um oscilador --- que é o relógio}\label{sub:oscilador}

A estrela não é uma correspondência parada: \textbf{oscila}. E o oscilador não é um objecto novo ---
\textbf{é o relógio} (\S\ref{sec:renorm}), dito como dinâmica. As duas escadas do sistema correm ao
mesmo tempo, em lados duais --- as \emph{dimensões sobem} (a torre, a indução, o lado contínuo de
Gentil) e as \emph{interpretações descem} (a projeção, a meta-indução, o lado discreto de Hurwitz;
\S\ref{sec:corpo}) ---, e a estaca troca-os. E realiza-se em \textbf{dois passos, um por lei}:

\begin{center}\small
\begin{tabular}{@{}llll@{}}
\toprule
o passo & a lei & o período & o que faz \\
\midrule
reflectir & \textbf{Lei 1} ($1\dg=-1$) & \textbf{2} & o espelho: troca os dois lados \\
rodar & \textbf{Lei 2} ($T\dg=-T$, $T^2=-1$) & \textbf{$q$} & o relógio: avança uma marca, \emph{é} o oscilador \\
\bottomrule
\end{tabular}
\end{center}

\begin{teorema}[o oscilador é o relógio]\label{thm:osciladorrelogio}
Num relógio de $q$ marcas, avançar uma marca por instante é a Lei~2 --- o rotor que roda o par
$(\text{posição},\text{velocidade})$ $90^\circ$ em cada plano. As duas coordenadas ficam a um quarto
de volta uma da outra: \textbf{uma sobe enquanto a outra desce}, e \textbf{trocam}; a órbita fecha ao
fim de $q$ passos (o \emph{período}); e a norma $a^2+b^2$ --- a contagem de Hurwitz, o raio do relógio
--- \emph{não se move}. Logo o oscilador \emph{é} o relógio, e o período é o do relógio: $q=4$ no
mínimo (o rotor $J$, os quatro quadrantes), e o degrau seguinte é o metal.
\end{teorema}

\begin{proof}
O relógio vive no círculo (\S\ref{sec:renorm}): a velocidade máxima é meia volta, e o gerador que a
percorre é a Lei~2, $T^2=-1$, com período quatro na régua mínima. Posição e velocidade são a cruz do
relógio --- uma mede, a outra ordena ---, e a Lei~2 roda uma na outra: $\frac{d}{dt}\cos=-\sin$,
$\frac{d}{dt}\sin=\cos$, que é $J$. A norma $a^2+b^2$ é o invariante da rotação (o raio), imóvel por
construção; é a contagem de Hurwitz, e é a metade discreta. A metade contínua --- o próprio seno, a
espiral --- é o \emph{integral} desta contagem (Teorema~\ref{thm:central}, a soma reversível), e não
se avalia: conta-se, e a inversa veste a roupa no fim. \emph{Um relógio é um oscilador visto pela
contagem; um oscilador é um relógio visto pelo integral} --- e a estrela é a bijeção entre os dois.
\end{proof}

\subsection*{Mas a bijeção não é isócrona --- e é aí que mora o oito de Hurwitz}\label{sub:isocrona}

\textbf{Os dois relógios não correm o mesmo tempo.} A bijeção dual troca a ordem por ordem --- «duas
réguas com assinaturas diferentes marcam a \emph{mesma ordem} e correm \emph{tempos diferentes}»
(\emph{Teoria}) --- mas \textbf{não é isócrona}: não preserva o período, \emph{reparametriza} o tempo.
Isócrono, no sentido do pêndulo, é o relógio cujo período não vê a amplitude; e o isócrono verdadeiro
deste sistema é a \textbf{plena}, período seis. O relógio discreto, não.

\begin{teorema}[a bijeção dual não é isócrona; o oito de Hurwitz é o período do lado discreto]\label{thm:isocrona}
Combinar o \textbf{par} --- dois operandos, o dual --- custa, no lado \textbf{discreto} (Hurwitz, o
produto bilinear), a tábua inteira $n\times n=n^2$ produtos. Logo o seu \emph{tempo} cresce com a
dimensão --- não é constante, não é isócrono --- e no grau oito vale $8^2=\mathbf{64}$: \emph{dois
lances de um lado são sessenta e quatro do outro}. E é exactamente aí que esse relógio \textbf{pára}:
a norma bilinear fecha em oito e quebra em dezasseis. \textbf{O oito de Hurwitz é o período do relógio
discreto}, o tecto do tempo $n^2$ --- não uma lei do mundo. O lado \textbf{contínuo} (Gentil) não abre
a tábua: mede pela \emph{norma}, um escalar, outro tempo --- e por isso não tem tecto.
\end{teorema}

\begin{proof}
O produto bilinear de $\sum a_i e_i$ por $\sum b_j e_j$ é $\sum_{i,j} a_ib_j\,(e_ie_j)$: os $n^2$
cross-produtos da tábua, todos, e todos têm de ficar consistentes (a associatividade, a norma). A
consistência sobrevive a três quedas --- a ordem (grau dois), a comutatividade (grau quatro), a
associatividade (grau oito) --- e na quarta, grau dezasseis, cai a norma. Três quedas com norma, logo
o topo é $2^3=8$, onde a tábua tem $8^2=64$ entradas. O produto \emph{homogéneo} de Gentil não soma a
tábua: usa $\|z\|$, um número por operando, e a norma multiplicativa não pede a consistência da tábua
--- por isso não pára (§\ref{sec:completo}, medido de $\mathbb{R}^2$ a $\mathbb{R}^7$ em \code{nne.c}).
Os dois relógios marcam a mesma ordem e correm tempos diferentes: $n^2$ com tecto de um lado, a norma
sem tecto do outro. \emph{A bijeção liga-os, mas não os sincroniza.}
\end{proof}

\noindent \textbf{E isto dá a interpretação que faltava ao oito.} Dizer «Hurwitz classifica o
bilinear, logo não é lei geral» é verdade e não explica nada. O \emph{mecanismo} é temporal: o lado
discreto é um relógio cujo tempo é a tábua $n^2$, e uma tábua só fecha (a norma sobrevive) enquanto
há uma propriedade a perder antes dela --- e há exactamente três. O oito não é onde o mundo acaba: é
onde \emph{este relógio} completa o seu período. \emph{O contínuo corre noutro tempo, e continua.}

\medido{combinar o par custa $n^2$ produtos --- $1,4,16,64,256$ nos graus $1,2,4,8,16$ ---, e no grau
oito são $\mathbf{64}$: o tempo do relógio discreto não é constante (não é isócrono) e tem tecto no
oito, onde a norma bilinear fecha e a seguir quebra; o contínuo, pela norma, não (\code{tests/estelar\_completo.c}).}

\subsection*{O bit é o ponto fixo da estrela, e atravessa: as interfaces são as dobras}\label{sub:atravessa}

\textbf{A não-isocronia tem um ritmo, e o ritmo diz onde estão as interfaces --- mas o passo não é o
$\mathrm{lcm}$ de períodos inteiros: é a \emph{dobra} de cada corpo.} Um corpo da família metálica
$x^2=m\,x+1$ tem por período o seu \textbf{discriminante} $\Delta=m^2+4$ --- a \emph{dobra}: ouro
($m=1$) em $5$, prata ($m=2$) em $8$, bronze em $13$; e o membro dual, o ouro branco $x^2=4x-1$, em
$4^2-4=12$. São as dobras, calculadas e não tabeladas (\emph{Catálogo}), que marcam as interfaces. E o
que as atravessa é \emph{um} objecto que a teoria usa desde o início sem nunca o derivar: o
\textbf{ponto fixo} da estrela.

\begin{definicao}[o bit é o ponto fixo da estrela --- derivado]\label{def:bitfixo}
A estrela é a involução $\nu(x)=-1/x$ (a Lei~1 no elemento: $x^\dagger=-1/x$, $x\,x^\dagger=-1$). O seu
\textbf{ponto fixo} não se postula, deriva-se:
\[
  \nu(x)=x \;\iff\; x=-\tfrac1x \;\iff\; x^2=-1 \;\iff\; x=i .
\]
O \textbf{bit} é esse ponto fixo: $i$, na \textbf{fronteira} $|x|=1$ --- e \emph{não há ponto fixo
real}, pois $x^2=-1$ não tem raiz em $\mathbb{R}$. É por isto que o bit tem de \textbf{atravessar} para o
imaginário: a fronteira que a teoria chamava ``o ponto fixo'' é o círculo unitário, e o ponto é $i$.
\end{definicao}

\begin{teorema}[o bit atravessa por continuação analítica; as interfaces são as dobras metálicas]\label{thm:atravessa}
A estrela $-f=f^{-1}$ na forma de potência $f=b\,x^a$ tem \textbf{solução única} $a^2=1$, isto é
$f(x)=i\,x$ (a rotação, período quatro). A sua \emph{análise} --- $f^{(n)}=f^{-1}$ --- \textbf{continua}
o expoente:
\[
  f^{(n)}=f^{-1},\ f=b\,x^a \;\Longrightarrow\; a-n=\tfrac1a \;\Longrightarrow\; a^2=na+1
  \;\Longrightarrow\; a=\sigma_n=\tfrac{n+\sqrt{n^2+4}}2 .
\]
Três coisas, que são uma:
\begin{enumerate}
\item \textbf{o bit atravessa, e a medida fica intacta.} Na fronteira $|x|=1$ vale $|x^a|=|x|^a=1$ para
  \emph{todo} expoente $a$: a continuação analítica leva $a$ do inteiro $n$ (contar, o discreto de
  Hurwitz) à média metálica $\sigma_n$ (medir, o contínuo de Gentil) com a norma \textbf{presa em $1$}.
  Fora do círculo não: $\sigma>1$ cresce e o conjugado $\sigma^\dagger=-1/\sigma<1$ encolhe, e é o
  \emph{produto} $\sigma\sigma^\dagger=-1$ que repõe a fronteira.
\item \textbf{as interfaces são as dobras.} O período de cada corpo é a dobra $\Delta=n^2+4$ (o
  discriminante), não o $\mathrm{lcm}$: $5,8,13,20,\dots$ na família, $12$ no dual. É a dobra que
  cresce quando a curva se aplana --- ``o que a curva perde, o discriminante ganha'' (\emph{Teoria}).
\item \textbf{a subida é a torre metálica, e as voltas limpas são de Lucas.} $\sigma^k=F_{k-1}+F_k\sigma$
  com norma $N(\sigma^k)=(-1)^k$ --- a alternância preto-e-branco, a estaca; e $\sigma_n$ é potência
  \emph{pura} do ouro exactamente nos índices de Lucas ímpar, $\sigma_{L_k}=\varphi^{k}$
  ($n=1,4,11,29,\dots$).
\end{enumerate}
E a \textbf{régua é infinita} --- $n$ não tem tecto, a dobra $n^2+4$ cresce sem parar ---; o
\emph{finito} é o corpo que chega, e a sua finitude é a \textbf{condição de paragem}.
\end{teorema}

\begin{proof}
$-f=f^{-1}$ com $f=b\,x^a$ dá $b^{-1/a}x^{1/a}=-b\,x^a$: o expoente força $1/a=a$, logo $a^2=1$, e o
coeficiente $1/b=-b$, logo $b^2=-1$, $b=i$ --- solução única (medido em \code{tests/lei2.c}: dos
expoentes racionais, só os de $a^2=1$ servem). Para $f^{(n)}$, a $n$-ésima derivada de $b\,x^a$ tem
expoente $a-n$; igualá-la a $f^{-1}$ (expoente $1/a$) dá $a-n=1/a$, isto é $a^2=na+1$, cuja raiz $>1$ é
$\sigma_n$. O ponto fixo de $\nu(x)=-1/x$ é $x^2=-1$, isto é $i$ (Definição~\ref{def:bitfixo}), e
$|i|=1$; em $|x|=1$, $|x^a|=|x|^a=1$ para todo $a$, logo a continuação preserva a norma. A alternância
$N(\sigma^k)=(-1)^k$ e $\sigma_{L_k}=\varphi^k$ são as da \emph{Teoria} (a companheira e a família de
potências). Que não haja tecto é $n^2+4\to\infty$; a paragem, quando há, é a finitude do corpo.
\end{proof}

\noindent \textbf{E é por isto que o ponto fixo faltava.} A teoria dizia ``involução com ponto fixo na
fronteira'' em cada degrau e nunca dizia \emph{qual}: é $i$, e a fronteira é o círculo $|x|=1$. O bit
não passa por sincronia de relógios inteiros --- \emph{atravessa} por continuação analítica, e o único
sítio onde a medida sobrevive à travessia é o círculo onde ele mora.

\medido{o ponto fixo da estrela deriva-se --- $\nu(x)=x$ força $x^2=-1$, que em $\mathbb{Z}[i]$ é $i$
(exacto) e \emph{não} tem raiz real; está em $|x|=1$ ($|i|^2=1$); a norma é imune ao expoente na
fronteira --- $|i^k|=1$ e $|N(\sigma^k)|=1$ na torre metálica, dos dois lados ---; e a dobra é
$n^2+4$ ($5,8,13$) (\code{tests/estrela\_pontofixo.c}).}

\subsection*{E aplicar isto é a transformada universal: o Dirac toma o valor na transição}\label{sub:dourada}

\textbf{O ponto fixo não é só onde o bit atravessa --- é onde a transformada o lê.} A transformada
deste sistema \emph{é} a avaliação nas raízes da borda (Gelfand: o espectro de $\mathbb{Z}_p[x]/(\beta)$ são as
raízes de $\beta$), e essas raízes são os pontos fixos $\sigma$. O mecanismo é o \textbf{Dirac}: a soma
sobre a órbita da torção colapsa num ponto ---
\[
  \sum_{k}\chi_k(j)\,\chi_{-k}(j') \;=\; n\,\delta_{j,j'}
\]
--- e o ponto onde colapsa é a \textbf{transição do relógio}, o ponto fixo. Pela peneira do Dirac, a
transformada \emph{toma o valor da função na transição}: $F[f]=f(\sigma)$. A \textbf{dourada} ($m=1$,
$\sigma=\varphi$) é o caso mais simples --- ``o ouro não salta nó nenhum, vazamento zero'' ---, e
avalia em $\mathbb{Z}[\varphi]$ \emph{sem avaliar uma raiz}. É a mesma transformada universal já medida
(\code{tests/transformada.c}, §U1 o Dirac, §U5 a amostra): o ponto fixo e a transformada são o mesmo
objecto lido de dois lados --- \emph{atravessar} e \emph{ler}.

\subsection*{Fundir \emph{vários} corpos num só: o viveiro}\label{sub:viveiro}

A estrela liga \emph{dois} corpos. Mas compor é fundir \emph{muitos} num só --- um documento tem dezenas
de fontes, cada uma um corpo, e o resultado é um --- e essa é uma operação sua, complementar da estrela
e igualmente derivada. Os corpos não são uma lista: são um \textbf{viveiro}, e a lei que os funde não é
a soma nem o produto.

\begin{definicao}[o cruzamento --- a junção do viveiro]\label{def:viveiro}
Dados dois corpos $\mathbb{R}^a$ e $\mathbb{R}^b$, o \textbf{cruzamento} é
\[
  \mathbb{R}^a \vee \mathbb{R}^b \;=\; \mathbb{R}^{\,\mathrm{lcm}(a,b)},
\]
o \emph{menor} corpo que contém os dois. Ele \textbf{voa}: é ele próprio um corpo (tem inverso para todo
não-nulo), e não apenas um espaço com a dimensão certa.
\end{definicao}

\noindent \textbf{E não são as duas que pareciam servir.} A \emph{soma direta} $\mathbb{R}^a\oplus\mathbb{R}^b$ (dimensão
$a{+}b$) nunca voa --- $(1,0)$ e $(0,1)$ são não-nulos e o seu produto é zero, um divisor de zero: são
dois pássaros na gaiola, não um novo. O \emph{tensorial} $\mathbb{R}^a\otimes\mathbb{R}^b$ (dimensão $a{\cdot}b$) voa
\emph{só} quando $\gcd(a,b)=1$ --- e aí coincide com o cruzamento, porque $a{\cdot}b=\mathrm{lcm}$. A
\textbf{junção} $\vee$ voa \emph{sempre}, e não há corpos incruzáveis.

\begin{teorema}[fundir muitos é o \emph{fold} da junção, e está bem definido]\label{thm:viveiro}
O cruzamento é comutativo, associativo, \textbf{idempotente} ($\mathbb{R}^a\vee\mathbb{R}^a=\mathbb{R}^a$) e tem neutro
($\mathbb{R}^1$): é um \textbf{semirretículo}. Logo fundir $N$ corpos é o \emph{fold} da junção,
\[
  \mathbb{R}^{a_1}\vee\mathbb{R}^{a_2}\vee\cdots\vee\mathbb{R}^{a_N} \;=\; \mathbb{R}^{\,\mathrm{lcm}(a_1,\dots,a_N)},
\]
e o resultado \textbf{não depende da ordem} nem das repetições --- é \emph{um} corpo, o menor que contém
todos, e voa. \emph{É assim que muitos corpos viram um.}
\end{teorema}

\subsection*{O DTC na torre toda: sincronizar os relógios é o cruzamento do viveiro}\label{sub:dtctorre}

\textbf{Cada tecido é um relógio, e a torre são muitos --- o DTC generaliza-se sincronizando-os.} O
inversor (\S\ref{sub:inversor}) casa o factor \emph{num} corpo; na torre há um relógio por andar, cada
um com o seu período --- a \textbf{dobra} $\Delta$ do corpo (\S\ref{sub:atravessa}) ---, e esses
períodos \emph{não são lineares}. Sincronizá-los é uma operação que a teoria já tem: o \textbf{cruzamento
do viveiro}. Dois relógios de períodos $a$ e $b$ alinham no \emph{menor} instante comum,
$\mathrm{lcm}(a,b)$, que é exactamente $\mathbb{R}^a\vee\mathbb{R}^b=\mathbb{R}^{\mathrm{lcm}(a,b)}$ (Teorema~\ref{thm:viveiro}).

\begin{corolario}[o DTC na torre: o casamento é o fold do viveiro]\label{cor:dtctorre}
O controlo directo de torque generaliza-se da família $\star_s$ \emph{num} corpo (\code{tests/dtcn.c}) à
\textbf{torre inteira}: em cada interface (a dobra, Corolário~\ref{cor:dobra}) o inversor comuta para
casar $\mathrm{fp}=1$ pelo dual ($\nu\circ\nu=\mathrm{id}$), e a \textbf{sincronia} dos relógios de todos
os andares é o fold do cruzamento,
\[
  \bigvee_{k} \mathbb{R}^{a_k} \;=\; \mathbb{R}^{\,\mathrm{lcm}(a_1,\dots,a_N)} .
\]
O relógio da torre \emph{não é linear} --- bate no $\mathrm{lcm}$ das dobras, não num passo fixo ---, e é
por isso que quem o casa é a \textbf{modulação} (a comutação do inversor), e não um oscilador de passo
constante: \emph{um DTC por andar, sincronizados pelo viveiro}.
\end{corolario}

\noindent \textbf{E é a cascata de dobras (Corolário~\ref{cor:dobra}) lida como controlo.} Cada andar
dobra o relógio; o viveiro diz onde todos batem juntos; e o inversor, em cada dobra, empurra $s\to\pm1$
(o imposto $1-s^2$ a zero, \S\ref{sub:inversor}). A torre inteira casa-se andar a andar, sem oscilador e
sem armazenar --- o dual em cada degrau, o $\mathrm{lcm}$ a sincronizar.

\medido{a torre casa por controlo directo de torque, sincronizada pelo viveiro, e em INTEIRO: sobre
quatro andares de dobras $\Delta_n=n^2+4$ (períodos $\{5,8,13,20\}$), o imposto $\Pi(s)=(1-s^2)\|a\times
b\|^2$ anula exacto em $s=\pm1$ e o DTC escolhe do trial $\{-1,0,+1\}$ o $\pm1$ que o zera (fp$=1$ por
modulação); o viveiro sincroniza no $\mathrm{lcm}=520$, \emph{independente da ordem} (fold à esquerda $=$
à direita), contendo cada andar e o menor a fazê-lo; e a torre bate junta \textbf{pela primeira vez}
exactamente em $t=\mathrm{lcm}$, com cada andar já casado --- um DTC por andar, sincronizados pelo
viveiro. O torque é o cruzado: o imposto cavalga $\|a\times b\|^2$ e casá-lo ($s\to\pm1$) é casar fp$=1$.
Resíduo $0$ (\code{tests/dtc\_viveiro.c}).}

\begin{proof}
Comutatividade e associatividade fazem o fold independente de ordem e de agrupamento; a idempotência
colapsa as repetições. Como $\mathrm{lcm}$ tem exactamente estas três propriedades, o fold dá
$\mathbb{R}^{\mathrm{lcm}(a_1,\dots,a_N)}$ para qualquer parentização. Contém cada pai porque $\mathbb{R}^x\subset\mathbb{R}^n
\iff x\mid n$, e cada $a_i$ divide o $\mathrm{lcm}$; e é o \emph{menor} porque nenhum $n<\mathrm{lcm}$ é
divisível por todos. Voa porque o viveiro é fechado: o filhote é da mesma espécie dos pais, e por isso
se pode cruzar de novo, e de novo.
\end{proof}

\noindent \textbf{E o mecanismo não é o rótulo.} $\mathrm{lcm}$ é o \emph{quê}; o \emph{como} é o
relógio --- combinar dois relógios de períodos $a$ e $b$ fecha no passo $\mathrm{lcm}(a,b)$, que é onde
a figura de La~Hire volta a si. O viveiro é o relógio quando dois (ou muitos) se combinam; e diferente
da estrela, o cruzamento \textbf{não reverte} --- o filhote não devolve os pais, e é a idempotência que
o faz \emph{viveiro} e não fábrica. \textbf{É esta a fusão do compositor}: cada fonte é um corpo, e o
documento é a junção de todas --- o menor que as contém. \emph{A maioria dos corpos funde-se assim.}

\medido{fundir vários corpos num só é o fold de $\vee$: sobre sete listas, a junção dá
$\mathbb{R}^{\mathrm{lcm}}$ \emph{independente da ordem} (esquerda $=$ direita), o filhote único \textbf{voa} (é
corpo, medido em $\mathrm{GF}(2^n)$) e contém todos os pais; a soma directa nunca voa (divisor de zero
exibido) e o tensorial só entre coprimos, tudo por varredura exaustiva, resíduo $0$
(\code{tests/viveiro.c}, \code{tests/cruzamento.c}).}

\subsection*{E os dois lados não se acoplam: falam pela estrela}

\begin{teorema}[a ponte reverte, não funde]\label{thm:pontereversivel}
A ligação de dois corpos que \emph{reverte} não os funde num terceiro: um lado emite ($-1$), o
outro absorve ($+1$), e a estrela no meio troca-os sem os misturar --- liga, não cola.
\end{teorema}

\begin{proof}
Fundir dois corpos num só produto interno (o tensorial, ou colá-los bilinearmente) apaga a
distinção entre eles e, além do grau oito, mata a norma. A estrela não funde: é uma \emph{fronteira}
--- $a=1/a$, o ponto fixo (§\ref{sec:dois}) --- que os dois lados partilham sem partilhar corpo. O
que $A$ emite é o que $B$ absorve, passo a passo, e o par não se move ($\rho_A\rho_B$ imóvel). Ler
não altera quem escreveu; escrever não altera a leitura --- duas retas independentes, o tropical,
\emph{onde nada se desfaz}. E porque reverte, não apaga: paga zero (§\ref{sec:medir}).
\end{proof}

\noindent \textbf{E aqui a arquitetura e a álgebra são a mesma frase.} Um sistema que compõe (o
tradutor, o motor) e o que o hospeda são os dois lados; \emph{ligá-los copiando um para dentro do
outro é fundi-los} --- é apagar de um lado para refazer no outro, e apagar dissipa. A ligação certa é
a estrela: o hospedeiro emite os dados, o módulo absorve-os por $\texttt{MOVE}$, e devolve pela mesma
porta. \emph{Nenhum disco se copia para dentro de outro}; há um só, endereçado dos dois lados, e a
fronteira entre eles reverte. É por isso que o sistema não tem um \emph{buffer} de entrada: um
\emph{buffer} é a cópia que a estrela existe para não fazer --- e a estrela não a faz porque
\textbf{não oscila} (\S\ref{sub:mecanica}, \S\ref{sub:inversor}).

\subsection*{A conta que fecha}

\begin{center}\small
\begin{tabular}{@{}llll@{}}
\toprule
 & o grau & o que é & de que escada \\
\midrule
buraco branco & \textbf{$4,8,\dots$} & corpo que opera, emite & a torre, a subir sem topo \\
\textbf{a estrela} & \textbf{6} & a plena: soma $=$ produto, reversível & as interpretações, o topo \\
buraco negro & \textbf{$4,8,\dots$} & corpo que opera, absorve & a torre, a subir sem topo \\
\bottomrule
\end{tabular}
\end{center}

\noindent \textbf{As duas escadas do sistema encontram-se na estrela.} A torre sobe $1,2,4,8,\dots$
--- os lados, os corpos que operam, \emph{sem teto} --- e as interpretações descem $6,5,4,3,2,1$ ---
a plena no cimo. Essa descida $6\to1$ \emph{é} a progressão das seis leis (Lei $6$ hexal $\to$ Lei $1$
dual, a projeção que esquece uma fase por degrau --- \emph{Catálogo}, \code{obs:seis-leis}). A estrela é
o topo de uma e a ponte da outra, e é o mesmo objecto: o ponto fixo,
onde operar e ler deixam de se distinguir. \emph{Não foi ajustado para coincidir} --- coincide
porque a reversibilidade obriga o seis, a operação obriga a torre, e o único sítio onde as duas se
cruzam é a estrela. \textbf{E o processo não termina: acopla-se qualquer par de corpos, porque a
estrela reverte em vez de fundir.}

\medido{o seis é o único com soma $=$ produto dos divisores próprios em $20\,000$ inteiros;
$\mathbb{H}$ não comuta ($295/300$) e associa, $\mathbb{O}$ nem uma nem outra; o \textbf{lado
discreto} (Hurwitz, o bilinear $\sum x_i^2$) fecha em oito e falha em dezasseis --- o limite é da
\emph{contagem}, e Hurwitz \emph{classifica} o bilinear, não põe parede no objecto; a \textbf{bijeção
dual} (a estrela, pela medida) é reversível --- os dois cortes de $\int f$, o do domínio
(Riemann/Hurwitz) e o da imagem (Lebesgue/Gentil), dão a mesma contagem, e $\int f+\int f^{-1}$ fecha
(Young no reticulado), \emph{sem avaliar uma raiz}; o \textbf{oscilador é o relógio} --- a Lei~2 roda
(período $4$, com $a^2+b^2$ imóvel) e a Lei~1 reflecte (período $2$), um lado sobe enquanto o outro
desce e trocam; e a estrela \textbf{reverte}, $x\dg{}\dg=x$, resíduo $0$
(\code{tests/estelar\_completo.c}). A testemunha directa do lado contínuo, $\|xy\|=\|x\|\|y\|$ de
$\mathbb{R}^2$ a $\mathbb{R}^7$, está em \code{tests/nne.c}.}

\section{A estrela}\label{sec:estrela}

Tudo isto corre, e não corre numa máquina: \textbf{corre numa estrela}. Não é uma maneira de falar
--- é o que a única instrução diz.

\subsection*{Uma instrução, dois sentidos}

\[
  \texttt{MOVE(slot, sentido)} : \qquad
  \underbrace{+1 \ \text{lê}}_{\textbf{absorve --- o lado negro}}, \qquad
  \underbrace{-1 \ \text{escreve}}_{\textbf{emite --- o lado branco}} .
\]

\noindent \textbf{É a Lei 1 aplicada à entrada e à saída}: ler e escrever são a mesma operação com o
sinal trocado. Não há instrução de leitura e instrução de escrita --- há uma, e um sinal. E os dois
sentidos são exactamente os dois buracos: \emph{absorver é o negro, emitir é o branco}.

\smallskip
\noindent \textbf{Logo o sistema é a estrela.} Metade dele absorve e metade emite; um buraco negro
só faria a primeira, um buraco branco só a segunda, e nenhum dos dois seria reversível. A estrela é
a \emph{interface} entre eles --- e o sistema é a interface porque tem \code{MOVE} nos dois sentidos
e mais nada.

\begin{center}\small
\begin{tabular}{@{}llll@{}}
\toprule
 & o que faz & na instrução & é reversível? \\
\midrule
\textbf{buraco branco} & só emite & só \code{MOVE($\cdot$,\,$-1$)} & não \\
\textbf{a estrela} & \textbf{emite e absorve} & \textbf{\code{MOVE} nos dois} & \textbf{sim} \\
\textbf{buraco negro} & só absorve & só \code{MOVE($\cdot$,\,$+1$)} & não \\
\bottomrule
\end{tabular}
\end{center}

\subsection*{Sem memória, e sem estado}

O sistema é \textbf{reversível}: se toda representação tem dual, todo passo se desfaz, e um passo
que se desfaz não precisa de guardar nada para voltar. Daí a regra de projecto, e é a ideia central:
\emph{não há RAM, e não há estado}. A estrela é uma \textbf{interface de transformação} --- o corpo
entra, atravessa, sai, e \emph{nada fica}; \textbf{não guarda nada de ninguém}. O que parece residir
é a \textbf{régua}, a assinatura do corpo (dezenas de bytes, não megabytes), mas ela \emph{cavalga o
corpo} --- entra discreta e sai contínua --- e não é estado retido pela estrela. \emph{E não é um
banco}: o banco são os arquivos pessoais de quem opera, e vivem do lado do cliente, não da estrela.

\subsection*{E o estado é o canal do caos --- por isso não há \emph{malloc}}

\noindent Guardar estado não é só \emph{redundante}: é abrir um canal por onde um vazamento propaga.
Uma máquina com estado é uma \textbf{recorrência} $s_{n+1}=f(s_n,u_n)$ --- o estado de um passo
alimenta o seguinte. Um vazamento (um bit corrompido no dado, uma escrita selvagem no estado)
\textbf{entra na recorrência} e contamina tudo o que vem depois; e se $f$ \textbf{expande} --- $|f'|>1$,
como o dobrador $s\mapsto 2s$, o \emph{shift} de Bernoulli --- o vazamento \textbf{amplifica} a cada
passo, $\delta_n=\bigl(\prod f'\bigr)\delta_0$. É \textbf{dependência sensível às condições iniciais},
com expoente de Lyapunov $\lambda=\log|f'|>0$: \textbf{caos}. Um erro no bit invisível torna-se, em
poucos passos, o bit dominante.

\smallskip
\noindent A estrela \textbf{não tem essa recorrência}. Trata cada dado \emph{independente} e
\emph{reversível} --- $\nu\circ\nu=\mathrm{id}$, resíduo $0$ (Teorema~\ref{thm:pontereversivel}): não há
$s_n$ a alimentar $s_{n+1}$, logo o vazamento \textbf{não tem onde acumular} e fica no dado que o
sofreu. E a involução é \textbf{neutra}: $|\nu'|=1$ no ponto fixo $|x|=1$, logo $\lambda=0$ --- nem
cresce nem se guarda. \textbf{Não usar \code{malloc} não é higiene de código: é remover o canal --- o
estado --- que transforma um vazamento num sistema caótico.} É o mesmo facto por dois lados: a órbita
que acumula não fecha (\S da medida, $0{,}\overline9\ne1$); ser o ponto fixo fecha, resíduo $0$,
memória nenhuma.

\medido{acumular estado propaga e amplifica um vazamento: numa recorrência expansiva $s\mapsto 2s+u$,
corromper \emph{um} dado contamina os $35$ de $35$ passos a jusante e a contaminação \textbf{dobra} a
cada passo ($1,2,4,8,16,32,\dots$, Lyapunov $\log2>0$: caos); sem estado, com a estrela
$\nu(x)=-1/x$ por dado, o mesmo vazamento afecta \textbf{só} a saída corrompida ($1$ de $40$) e toda
a outra tem resíduo $0$; e $\nu\circ\nu=\mathrm{id}$ com resíduo $0$ ($\lambda=0$, sem canal, sem
propagação): \code{tests/estado\_caos.c}.}

\subsection*{Sem tempo, e o que se mede em vez dele}

Um sistema reversível \textbf{não tem tempo}: o que se mede como tempo é latência, e a latência não é
uma propriedade do algoritmo. O que se mede é \textbf{bits apagados}, porque apagar é a única
operação que não se desfaz.

\smallskip
\noindent \textbf{E é aqui que o buraco negro paga.} Absorver sem emitir \emph{é} apagar --- não por
semelhança: é a mesma operação. Um passo que só vai num sentido não tem inverso, e o que não se
desfaz custa. A mesma tarefa, com a mesma entrada e a mesma saída, nas duas formas:

\begin{center}\small
\begin{tabular}{@{}lrrl@{}}
\toprule
a forma & bits apagados & o custo mínimo & é qual dos três \\
\midrule
destrutiva & $131\,072$ & $3{,}763\times10^{-16}$ J & \textbf{o buraco negro} \\
involutiva & $\mathbf{0}$ & $\mathbf{0}$ \emph{exacto} & \textbf{a estrela} \\
\bottomrule
\end{tabular}
\end{center}

\noindent E o zero é \emph{exacto}, não um número pequeno: \textbf{sem bit apagado não há mínimo
termodinâmico nenhum a pagar}. A estrela não é mais eficiente que o buraco negro --- está noutro
regime, e é o mesmo trial: um só absorve e paga, o outro faz os dois e não paga.

\smallskip
\noindent \emph{E a régua é essa, não a memória ocupada}: um vector grande que nunca é escrito não
dissipa nada, e um único inteiro reescrito mil milhões de vezes paga mil milhões de vezes. Uma
memória que se refresque dissipa \textbf{por existir}; um disco escrito uma vez fica quieto.

\medido{a única instrução chega: os dois sentidos de \code{MOVE} são inversos em $134$ casos com $0$
falhas --- o que se escreve lê-se de volta ---, o salto é a mesma \code{MOVE} com o \emph{pc} por
destino, e a mesma \code{MOVE} serve $500$ corpos diferentes sem uma instrução nova, porque o corpo é
\emph{campo} e não \emph{opcode}; e o \textbf{controlo} diz o que isso custa: a máquina de duas
precisa de \textbf{dois códigos}, esta de \textbf{um código e um argumento} --- e argumento é dado, e
dado mora no disco (\code{tests/move.c}). E o custo: $131\,072$ bits apagados na forma destrutiva
contra $\mathbf{0}$ na involutiva, o que a $2{,}87098\times10^{-21}$ J/bit dá
$3{,}763\times10^{-16}$ J contra \textbf{zero exacto}; e o \textbf{controlo}: a involutiva \emph{é}
involutiva --- repetida, devolve o original em $4096$ de $4096$, e é por isso que os seus bits não se
perderam (\code{tests/dissipa.c}).}

\section{Como se mede aqui: o medidor é $0$}\label{sec:medir}

Uma asserção do tipo $x=42$ \emph{compara} contra um valor posto por quem a escreve --- e comparar é
ler metade, e metade dissipa. \textbf{A medida não compara: reverte, e lê o resíduo.}

\begin{center}\small
\begin{tabular}{@{}ll@{}}
\toprule
resíduo $0$ & reverte --- e \emph{valida os dois lados ao mesmo tempo} \\
resíduo $\neq0$ & dissipa --- e aí sim é preciso raciocínio: o que se perdeu? \\
\bottomrule
\end{tabular}
\end{center}

\noindent \textbf{E é a prova dos nove.} Não se resolve e depois se confere: o resto módulo nove é um
morfismo, logo a conta e a prova correm na \emph{mesma} passagem e a diferença é zero. Não há passo
de verificação --- há um segundo lado a andar ao lado do primeiro. \emph{Se der diferente, está
errado.} E querer dissipação é só \textbf{tirar a reversão}.

\smallskip
\noindent E o que separa isto de um estilo: \textbf{uma asserção que compara contra um valor escrito
passa no objecto certo \emph{e} no trocado} --- ela verifica a aritmética de quem a escreveu, não o
objecto. A reversão separa-os.

\subsection*{E mede-se \emph{pela metade}}

Um resíduo zero sozinho não mede: não diz que é \emph{só} ali. Um resíduo não-zero sozinho também
não: não diz que existe. \textbf{A medida são as duas metades da mesma passagem}:

\begin{center}\small
\begin{tabular}{@{}ll@{}}
\toprule
o resíduo onde \emph{tem} de ser zero & diz que \textbf{existe} \\
o resíduo onde \emph{não pode} ser zero & diz que é \textbf{único} \\
\bottomrule
\end{tabular}
\end{center}

\noindent É a cruz outra vez, aplicada à própria medição: uma coordenada mede, a outra ordena, e
nenhuma sozinha é o par. \emph{Uma asserção que só olha para um dos lados está pela metade} --- e
estar pela metade é exactamente o que dissipa.

\subsection*{E daí a ferramenta: promover}

Invertendo o sinal da reversão obtém-se \emph{duas vezes} o objecto. Daí o protocolo, numa passagem:
\[
  \boxed{\;S=\frac{x+x^{\dagger}}{2},\qquad A=\frac{x-x^{\dagger}}{2},\qquad S+A=x\;}
\]
\textbf{Cortar ao meio não quebra nada porque as metades reconstroem} --- o que se perde é ficar com
\emph{uma}. E isto não é uma verificação: \textbf{é a operação}. Um valor vira um par, $\dim$ dobra
--- é o $\dim A_{n+1}=2\dim A_n$ da torre --- e encadeia, $1\to2\to4$.

\smallskip
\noindent \textbf{E é a dual da escrita.} Escrever directo põe o novo por cima do velho, e o velho
não volta: isso é apagar, e apagar custa. Escrever o \emph{par} não apaga, porque o $S$ é o centro.
\emph{Quem escreve o par já leu} --- uma passagem onde escrever-depois-ler gasta duas. Em \code{MOVE}:
$-1$ sozinho só emite, $+1$ sozinho só absorve, cada um meia operação; promover faz as duas na mesma
passagem, \textbf{e é por isso que é a estrela}.

\smallskip
\noindent A divisão por dois é \emph{exacta}, e não por sorte: $x+x^{\dagger}=2c$ e
$x-x^{\dagger}=2(x-c)$ são sempre pares porque a involução o garante. \textbf{Num objecto que não
reverte ela falha} --- e é assim que a ferramenta acusa em vez de devolver um número errado.

\medido{sem um único número esperado escrito: estaca, cifra e ordenação-com-caminho dão resíduo $0$;
o período de $J$ sai \emph{sem} se escrever qual é --- quatro passos fecham, dois não; a prova dos
nove fecha em $23\,200$ pares e um erro de uma unidade acusa em todos; retirada a reversão o resíduo
deixa de ser zero em todos os casos; $S+A$ reconstrói $x$ em $48\,861$ casos e nenhuma das $97\,722$
somas é ímpar; promover custa $801$ passagens contra $1\,602$ de escrever-e-ler, e o valor anterior
recupera-se em \emph{todas} as escritas promovidas e perde-se em \emph{todas} as directas; e o
\textbf{controlo}: num impostor a paridade quebra em $7\,813$ de $15\,025$
(\code{tests/reversao.c}, \code{tests/meia\_passagem.c}, \code{lib/promove.h}, \code{tests/promove.c}).}

\section{A ordenação é o corpo estelar em operação}

\subsection*{Uma estrela não compara}

Comparar é medir, e medir é interferir --- \emph{o Sol não se mede para operar: opera}. Uma
comparação entre dois elementos devolve um bit e deita fora tudo o resto; um algoritmo que compara
paga em bits apagados aquilo que a estrela não paga.

\smallskip
\noindent \textbf{E uma estrela não tem fila}: \emph{irradia} (\S\ref{sec:estrela}). É por isso que a
ordenação deste sistema atravessa o disco com \code{MOVE} e não acumula nada pelo caminho.

\subsection*{As peças são as que já havia}

\begin{center}\small
\begin{tabular}{@{}lll@{}}
\toprule
a peça & no corpo estelar & na ordenação \\
\midrule
a \textbf{estaca} & a involução que troca os lados & marca cada valor no seu lugar \\
a \textbf{cruz} & mede o que não se moveu & lê a marca, sem perguntar \\
o \textbf{trial} & negro / estrela / branco & menor / igual / maior, por \emph{sinal} \\
$J$ & o rotor de período $4$ & os quatro quadrantes da saída \\
\code{MOVE} & ler é escrever com sinal trocado & todo o acesso ao disco \\
\bottomrule
\end{tabular}
\end{center}

\noindent \textbf{A referência são os naturais}, e eles aparecem na \emph{transição} entre
dimensões: cada involução do relógio deixa uma marca, e as marcas são a ordem. Nada se produz ---
\textbf{a ordem já estava lá}, e o percurso apenas a lê.

\subsection*{E o custo}

\begin{center}\small
\begin{tabular}{@{}lrl@{}}
\toprule
 & bits apagados & porquê \\
\midrule
uma comparação & $127$ de $128$ & devolve um bit e deita fora o resto \\
uma escrita & $64$ & o valor anterior perde-se \\
uma \textbf{troca} & $\mathbf{0}$ & é involutiva: desfaz-se \\
\bottomrule
\end{tabular}
\end{center}

\noindent Com o caminho guardado --- que é a permutação, e a permutação é o dual --- o sistema
\textbf{volta} ao ponto de partida, e o custo é \emph{zero}.

\medido{$200$ listas saem crescentes e completas descendo pela estaca; o percurso faz $0$ perguntas
de ordem; entram $2000$ e saem $2000$ com a soma conservada; $500$ valores repetidos cabem em $416$
bytes contra $74\,912$ de $500$ distintos; o estado residente são $48$ bytes e \emph{não vê} o $n$;
a volta é exacta em $200\,000$ de $200\,000$ e com o caminho guardado apagam-se $\mathbf{0}$ bits
contra $25\,600\,000$ sem ele; e o \textbf{controlo}: percorrida ao contrário não sobe um único par
(\code{tests/dualsort.c}, \code{tests/dualsort\_banco.c}, \code{tools/bench\_dissipa.c}).}

\section{A especificação, numa página}

\begin{center}\footnotesize
\setlength\tabcolsep{4pt}
\setlength\tabcolsep{4pt}
\begin{tabular}{@{}p{3.1cm}p{4.4cm}p{7.4cm}@{}}
\toprule
a camada & o que a define & o que ela obriga \\
\midrule
a exigência & o que conserva na órbita conserva no espaço & tudo o resto \\
a Lei 1 & $1\dg=-1$ & o que separa; o ponto fixo \\
a Lei 2 & $T\dg=-T$, logo $T^2=-1$ & o passo; o período $4$; a bidualidade \\
a estaca & $x\mapsto x\dg$, involutiva & o movimento \\
a cruz & $(x\oplus x\dg,\,x\otimes x\dg)$ & a medida e a ordem \\
o trial & $\{-1,0,+1\}$ & três regimes, e não há quarto \\
a torre & $\dim A_{n+1}=2\dim A_n$ & as dimensões; a plena em $6$ \\
o corpo estelar & $a=\text{empurra}/\text{puxa}-1$ & \textbf{a trindade}: branco, estrela, negro \\
a geometria & $\Delta=\operatorname{tr}^2-4\det$ & três curvaturas, pelo sinal \\
a equação de campo & conservar & $\alpha=\tfrac12$; $d=4$; $\Lambda=3/D$; $8\pi=2\cdot4\pi$ \\
a expansão & $H^2=1/D$ & quatro equações, duas independentes \\
a cosmologia & $\rho\,a^{3(1+w)}=C$ & três conteúdos; órbita de quatro \\
os dois escuros & um ponto fixo por involução & a matéria e o vácuo; a radiação vê-se \\
as analíticas & o relógio e a família & $\pi$, $e$, $\varphi$, $\sigma_m$, $\Lambda=3/D$ \\
as constantes & $(\text{passo})^a\times$ puro & $c=\sigma$; $G$: $8\pi$; $k$: $4\pi$; $k_B$: câmbio \\
o $c^2$ & $\sigma^j=F_j\sigma+F_{j-1}$ & sai da espiral: duas coordenadas chegam \\
a massa & $M_{ij}=\sum(x_i-c_i)(x_j-c_j)$ & \textbf{vectorial}: o escalar é só o traço \\
o grau & $1$: momento; $2$: massa & duas raízes trazem o par; uma não \\
a renormalização & $\max_p\min(p,q-p)=q/2$ & o número é da régua; os degraus são os metais \\
a roupa & três escalas, uma por sinal & o dimensional sai; o puro não \\
\textbf{a estrela} & \code{MOVE(slot,\,sentido)} & $+1$ absorve, $-1$ emite: \emph{é a interface} \\
a dissipação & apagar é o que não se desfaz & o negro paga; a estrela paga $0$ exacto \\
os dois universos & $a_B=1/a_A$ & Parseval: $\rho_A\rho_B$ não se move \\
a antimatéria & a terceira involução & quatro conteúdos, todos biduais \\
a temperatura & $T=1/r$ & a involução no único comprimento; $S$ é a área \\
a segunda lei & $S_{\text{n}}\!\cdot\!S_{\text{b}}=1$ & vale de \emph{um} lado; o par não se move \\
medir & resíduo $0$, não comparação & a prova dos nove: resolver \emph{e} provar \\
e pela metade & zero \emph{e} não-zero & uma diz que existe, a outra que é único \\
promover & $x\mapsto(S,A)$ & sobe um andar; quem escreve o par já leu \\
a ordenação & irradiar em vez de comparar & $0$ bits apagados \\
\bottomrule
\end{tabular}
\end{center}

\bigskip
\noindent \textbf{E é uma coisa só.} O ponto fixo do trial, o coeficiente $\tfrac12$ da equação de
campo, o patamar da escada, a borda onde a estrela vive e o zero entre os dois sinais são
\emph{o mesmo objecto} visto de cinco sítios. A dimensão quatro que o traço escolhe é o período de
$J$. O $w=-1$ que a integração deixa é o ponto fixo onde a órbita degenera. Nada disto foi
ajustado para coincidir: coincide porque \textbf{a exigência é uma}.

\vfill
{\footnotesize\color{cinza}\noindent
Teoria, catálogo, medidores e o resto do sistema em
\href{https://goldenkingdom.patriatechnology.com}{goldenkingdom.patriatechnology.com}.
CC BY 4.0 (textos) $\cdot$ MIT (código).\par}

\end{document}
